% HistoCrypt 2027 regular paper draft (10 pages excluding references). Anonymous for submission.
% Compile: tectonic llm-ciphers.tex   (or pdflatex + bibtex + pdflatex x2)
% Every number in the \N... macros and in the tables is generated by papers/lasry/_analysis.py from the
% <target>/profile.json files; rerun it after any profile change, then compile.
\documentclass[11pt]{article}
\usepackage{histocrypt}
\usepackage{mathptmx}
\usepackage[utf8]{inputenc}
\usepackage[T1]{fontenc}
\usepackage{url}
\usepackage{latexsym}
\usepackage{amssymb}
\usepackage{booktabs}
\usepackage{longtable}
\usepackage{graphicx}
\special{papersize=210mm,297mm}
% generated by papers/lasry/_analysis.py; do not edit
\newcommand{\Ntargets}{315}
\newcommand{\Ndocs}{863}
\newcommand{\Nworked}{292}
\newcommand{\Nobtained}{233}
\newcommand{\PctObtained}{80}
\newcommand{\Ncomplete}{164}
\newcommand{\Npartial}{69}
\newcommand{\NCT}{29}
\newcommand{\NCTcomplete}{19}
\newcommand{\NCTpartial}{10}
\newcommand{\NEPT}{11}
\newcommand{\NEPTcomplete}{8}
\newcommand{\NEPTpartial}{3}
\newcommand{\NAPT}{39}
\newcommand{\NAPTcomplete}{19}
\newcommand{\NAPTpartial}{20}
\newcommand{\NEK}{17}
\newcommand{\NEKcomplete}{12}
\newcommand{\NEKpartial}{5}
\newcommand{\NKK}{81}
\newcommand{\NKKcomplete}{53}
\newcommand{\NKKpartial}{28}
\newcommand{\NDEC}{56}
\newcommand{\NDECcomplete}{53}
\newcommand{\NDECpartial}{3}
\newcommand{\NNS}{59}
\newcommand{\NNScomplete}{0}
\newcommand{\NNSpartial}{0}
\newcommand{\NNA}{23}
\newcommand{\NNAcomplete}{0}
\newcommand{\NNApartial}{0}
\newcommand{\Nrecovered}{79}
\newcommand{\Napplied}{98}
\newcommand{\Ndeciph}{56}
\newcommand{\Nrecoveredcomplete}{46}
\newcommand{\NnotCipher}{12}
\newcommand{\NoffLine}{8}
\newcommand{\NnaOther}{3}
\newcommand{\NnotSolvedByOthers}{2}
\newcommand{\YearMin}{1398}
\newcommand{\YearMax}{1945}
\newcommand{\Nundated}{12}
\newcommand{\TokensTotal}{854{,}881}
\newcommand{\Nlengthknown}{241}
\newcommand{\Nmeasured}{169}
\newcommand{\Nlanguages}{13}
\newcommand{\NsixteenthC}{148}
\newcommand{\NseventeenthC}{70}
\newcommand{\NfifteenthC}{13}
\newcommand{\NfourteenthC}{1}
\newcommand{\NeighteenthC}{49}
\newcommand{\NnineteenthC}{15}
\newcommand{\NtwentiethC}{7}
\newcommand{\NmodernC}{22}
\newcommand{\Ncto}{66}
\newcommand{\NctoCT}{23}
\newcommand{\NctoOther}{4}
\newcommand{\NctoDEC}{4}
\newcommand{\NCTnotcto}{6}
\newcommand{\NctoNot}{39}
\newcommand{\NctoPct}{35}
\newcommand{\Ncrib}{27}
\newcommand{\Ncribobtained}{25}
\newcommand{\Nkeyed}{113}
\newcommand{\Nkeyedobtained}{107}
\newcommand{\Nkeyedpct}{95}
\newcommand{\Ndecstart}{48}
\newcommand{\Ndecstartobtained}{47}
\newcommand{\Nclean}{82}
\newcommand{\Ncleanobtained}{44}
\newcommand{\Ncleanrecovered}{28}
\newcommand{\NcleanCT}{14}
\newcommand{\NcleanApplied}{16}
\newcommand{\Ncleancomplete}{29}
\newcommand{\NCTprior}{15}
\newcommand{\NCTclean}{14}
\newcommand{\NCTpriorbefore}{5}
\newcommand{\NpriorFoundDuring}{75}
\newcommand{\NpriorFoundBefore}{102}
\newcommand{\Nprintprior}{45}
\newcommand{\Nrecalpha}{24}
\newcommand{\Nrecrandom}{16}
\newcommand{\NCTalpha}{8}
\newcommand{\NCTrandom}{6}
\newcommand{\Nkeyorderunknown}{187}
\newcommand{\Nsteps}{3{,}376}
\newcommand{\Nfailedsteps}{810}
\newcommand{\PctFailedSteps}{24}
\newcommand{\Nsolversteps}{301}
\newcommand{\Nlitsteps}{435}
\newcommand{\Ntranssteps}{409}
\newcommand{\Nsolverfailed}{202}
\newcommand{\Ncontrol}{112}
\newcommand{\NnotControl}{13}
\newcommand{\NCTcontrol}{6}
\newcommand{\NCTverified}{26}
\newcommand{\NCTsolverstep}{24}
\newcommand{\Ndocsllm}{580}
\newcommand{\Ndocsdecode}{148}
\newcommand{\Ndocspublished}{61}
\newcommand{\PctDocsllm}{67}
\newcommand{\Nimg}{192}
\newcommand{\Nimgobtained}{150}
\newcommand{\Ndeciphtranscribed}{18}
\newcommand{\Ncatcorr}{138}
\newcommand{\Nsenderid}{120}
\newcommand{\Nkeyextended}{73}
\newcommand{\Nambiguous}{20}
\newcommand{\Nfrac}{173}
\newcommand{\Nfracmeasured}{103}
\newcommand{\Nbarup}{41}
\newcommand{\Nbardown}{10}
\newcommand{\Npartialhigh}{8}
\newcommand{\Nclaude}{289}
\newcommand{\Ngpt}{18}
\newcommand{\Nboth}{5}
\newcommand{\Nmodelunknown}{13}
\newcommand{\DateFirst}{14 September 2026}
\newcommand{\DateLast}{9 October 2026}
\newcommand{\Nhumaninput}{5}
\newcommand{\Nreview}{290}
\newcommand{\Nreviewagree}{216}
\newcommand{\PctReviewagree}{74}
\newcommand{\Nreviewflagged}{115}
\newcommand{\Nreviewchanged}{74}
\newcommand{\Nflaggedchanged}{37}
\newcommand{\Ncommits}{1{,}834}
\newcommand{\Ncommitsmodel}{1{,}760}
\newcommand{\Nexternal}{12}
\newcommand{\Nexternalworked}{12}
\newcommand{\Nexternalobtained}{10}
\newcommand{\Nnotes}{314}
\newcommand{\Ncode}{225}
\newcommand{\ContribTargets}{aerssen1601 (key from adjacent plaintext, partial); battier1686 (known key applied, partial); beuningen1676 (key from adjacent plaintext, partial); champagne1590 (key from adjacent plaintext, partial); dedem1788 (not solved); dinteville1592 (key from adjacent plaintext, complete); erving1807aug (key from external plaintext, complete); foix1563 (existing decipherment, partial); labbe1582 (key from ciphertext only, complete); malsburg1637 (key from ciphertext only, partial); matignon1586 (known key applied, partial); palm1727 (not solved)}
\newcommand{\TopLanguages}{French 97, Italian 56, English 55, Spanish 35, German 31, Latin 19}


% Submission version: anonymous. Camera-ready: compile llm-ciphers-camera.tex, which defines \cameraready and so
% names the authors and credits the second contributor by name (as the write-ups on the project site do) in the
% text, the tables and the acknowledgements.
\newif\ifanon
\ifdefined\cameraready \anonfalse \else \anontrue \fi
\newcommand{\Who}[2]{\ifanon #1\else #2\fi}
\newcommand{\ContribCredits}{George Lasry (77), Brice Dudout (50), WAINUT (30), Michael Brackx (10), Arya Sanketbhai Patel (8), Harumi-Matsuba (7), Feyseel Nur (6), alejandrozarco (2), Angus Robinson (2), Brett Burford (2), Larry Beck (2), Harrison Bornstein (1), Paolo Rosson (1), johonn (1), Satoru (1), setsunaatto (1), NoAutopilot (1)}


\title{Large Language Models against Historical Ciphers: \\ A Profiled Record of \Nalltargets\ Targets, \YearMin--\YearMax}

\ifanon
\author{Anonymous submission}
\else
\author{Daniel Bourdeau \\ Independent researcher \\ {\tt dnbourdeau@gmail.com} \And George Lasry}
\fi
\date{}

\begin{document}
\maketitle

\begin{abstract}
We report what agentic large language models (LLMs) achieved on historical cipher targets dated \YearMin--\YearMax, mostly diplomatic letters in archives, attacked between \DateFirst\ and \DateLast. The record has \Nalltargets\ targets: \Ntargets\ worked by the project, which the figures below describe, and \Ncontrib\ solved by \Ncontributors\ outside contributors, who used their own methods and often their own LLMs, and whose results are reported apart. Each target has a fixed profile: the cipher system, the documents, what the model started from, whether a reading already existed, the ordered solution steps, and an outcome split into \emph{method} (how the text was obtained) and \emph{extent} (complete or partial). The text was obtained for \Nobtained\ of the \Nworked\ targets that could be worked (\PctObtained\,\%), but that figure hides the methods. The key was recovered here for \Nrecovered: \NCT\ from the ciphertext alone, \NAPT\ from plaintext beside the cipher and \NEPT\ from plaintext found elsewhere. For \Napplied\ a key that already existed was found and applied. For \Ndeciph\ a decipherment already existed. From ciphertext alone the model recovered the key in \NctoCT\ of \Ncto\ attempts. On the \Nclean\ targets with no known earlier reading, the text was obtained for \Ncleanobtained. The model's main strength was locating keys, cribs and decipherments; its main failures were noisy transcriptions, short texts and large nomenclators without key material.
\end{abstract}

\section{Introduction}

Historical cryptology has been reshaped twice in the last two decades: by metaheuristic search, which made ciphertext-only attacks on homophonic and transposition ciphers routine \cite{Lasry:2018}, and by shared collections such as DECODE, which put thousands of ciphertexts and keys in one place \cite{Megyesi:2019,Megyesi:2020}. Large language models add a third element. An agentic model can read a manuscript image, search archive catalogues and printed editions, write and run a solver, and judge whether its output is language. Whether that adds up to useful cryptanalysis, and on which ciphers, is an empirical question.

This paper answers it from a working record rather than a benchmark. Between \DateFirst\ and \DateLast\ agentic models worked through lists of unsolved or unread historical ciphers: Tomokiyo's list of unsolved historical ciphers \cite{Tomokiyo:unsolved}, Schmeh's list of the top fifty unsolved encrypted messages \cite{Schmeh:top50}, the DECODE records marked non-decrypted, and a sweep of the manuscript catalogue of the Biblioth\`eque nationale de France. All of the work was done in one version-controlled repository, and that repository is the record we analyse: a folder per target holding the agent's working notes, the transcriptions, the solver code and the readings, and a commit history (\Ncommits\ commits over \Ndays\ days) that dates every step and names the model that made it (Section~\ref{sec:method}). The weakness of such a record is that each target's notes describe its cipher in their own words. To make the targets comparable, we defined a fixed profile and filled it for every target from the notes and the files, measuring counts wherever they could be measured.

A record of this kind also forces a question the field has left informal: what counts as a success. A letter can be ``deciphered'' because the key was rebuilt from nothing, because a key sheet was found in another volume, or because a contemporary clerk had already written the plaintext between the lines. We therefore classify every outcome by method, following the categories proposed by the second author, and report the extent of the text obtained separately.

Our contributions are:
\begin{itemize}
\item a profile schema for historical cipher challenges and their solutions, in which the factors that plausibly affect success are recorded the same way each time (Section~\ref{sec:schema});
\item profiles for \Ntargets\ targets worked by the project, with \Ndocs\ ciphertext documents from \YearMin\ to \YearMax\ in \Nlanguages\ plaintext languages (Section~\ref{sec:corpus}), and for \Ncontrib\ more targets solved by outside contributors, reported separately (Section~\ref{sec:contrib});
\item an analysis of outcome method and extent by starting point, prior reading, cipher family and length (Section~\ref{sec:results});
\item six case studies, one per outcome route, and a discussion of the failures (Sections~\ref{sec:cases}--\ref{sec:failures}).
\end{itemize}

\section{Related work}

Ciphertext-only attacks on historical ciphers by hill climbing and simulated annealing are reviewed by \newcite{Lasry:2018}, whose methodology covers homophonic, transposition and machine ciphers; recent HistoCrypt and \emph{Cryptologia} papers apply it to single nomenclators \cite{Lasry:2023} and to whole correspondences \cite{Lasry:2023mary}. The scoring function in such attacks is an $n$-gram language model, and \newcite{Megyesi:2023lm} show that models built from period texts improve decipherment of seventeenth-century and older homophonic ciphers. The DECODE database \cite{Megyesi:2019,Megyesi:2020} standardises the metadata of ciphertexts and keys (provenance, dating, symbol set, transcription); our profile borrows its document fields and adds fields about the attack and its result.

Neural models have been trained to decipher directly. Sequence-to-sequence models learn to crack simple substitution ciphers \cite{Aldarrab:2021}, and a causal language model trained on symbol recurrences recovers keys of real homophonic ciphers \cite{Kambhatla:2023}. Bayesian inference had earlier reached 97.8\,\% on the Zodiac-408 cipher \cite{Ravi:2011}. These systems are cipher solvers; they do not find sources or decide what a document is.

General-purpose LLMs have been evaluated on ciphers as reasoning tasks. CipherBank \cite{Li:2025} tests chat and reasoning models on nine encryption algorithms applied to synthetic plaintexts and finds that even the strongest models solve fewer than half of tasks that human cryptanalysts can; \newcite{Maskey:2025} benchmark LLMs on decrypting ciphertexts of several classical and modern schemes. Both use generated ciphertexts with known answers and no archival context. We know of no evaluation of LLMs on archival ciphertexts, where most of the work lies before and around the cryptanalysis: finding the document, transcribing it, identifying the system, and finding the key or a decipherment if one survives.

\section{The profile}
\label{sec:schema}

A profile is a JSON file in five parts, validated against a schema. Every field may take the value \emph{unknown}; a gap is then explicit and counted, rather than filled by a plausible guess.

\paragraph{System.} The cipher types that occur (monoalphabetic, homophonic, polyphonic, syllabic, nomenclator, code, transposition, polygraphic, polyalphabetic, superencipherment; several at once means a combined system). The symbol kind (digits, letters, graphic symbols or mixtures), digit-group width and separation, and the number of distinct symbols. Diacritics and what they do. Homophones per vowel and per consonant. The nomenclator: presence, size, kinds of entry and the number of code groups still open. Nulls, word division and doubled letters. \emph{Key order}: whether the key assigns its signs in alphabetical order, so that a few known values predict the rest. Any sibling key from the same office.

\paragraph{Documents.} One entry per ciphertext unit: shelfmark, date, issuing power, route, plaintext language, pages. Whether the document contains clear text, and whether a plaintext exists and where: interlinear, in the margin, on a separate folio or in a separate collection, in a printed edition, or in an unrelated source. The length in enciphered tokens and the number of distinct tokens, flagged as measured when a common script counted them from the transcription file. Who transcribed the text (the model from images, a human, a published or DECODE transcription) and the image quality.

\paragraph{Conditions.} Whether a reading of the text existed anywhere before this project obtained it, where, whether it was found before, during or after the attempt, and whether the work used it. This is the contamination question: a model may have seen a printed decipherment in training. The inputs the attack started from and the attack class (ciphertext-only, crib, known plaintext, archive decipherment, sibling key, key from archive, published key). The human's role, the tools, the model identifiers and the dates.

\paragraph{Solution.} The ordered steps, each with a kind (access, transcription, statistics, hypothesis, crib, solver, control, sibling key, key from source, prediction, reading, verification, literature search) and a result (worked, partial, failed, ruled out). Failed steps are recorded.

\paragraph{Outcome method.} How the text was obtained, in one of eight categories. Three recover a key here, by cryptanalysis: \emph{key recovered from ciphertext only} (frequencies, solvers, pattern words, the order of the key; no plaintext and no existing key); \emph{key recovered from adjacent plaintext} (a partial decipherment between the lines or in the margin, clear passages in the same letter, a deciphered copy in the same file); and \emph{key recovered from external plaintext} (a printed edition, a calendar of state papers, a copy in another archive). Two apply a key that already existed: \emph{external key matched} (a key for another correspondent or series, or a key sheet filed elsewhere, found to fit) and \emph{known key applied} (a key already tied to the document or its correspondence). One transcribes: \emph{existing decipherment} (a complete decipherment already existed; the work was transcribing the cipher and checking the decipherment against it). The last two are \emph{not solved} and \emph{not applicable} (not a cipher, or the item could not be obtained). Applying a known key is not mechanical when the key is polyphonic or variable-length, and the profile records such cases in a separate contribution field.

\paragraph{Extent.} How much of the text was obtained: complete or partial. The field has no agreed threshold. DECODE marks each record decrypted, partly decrypted or non-decrypted, with the status ``freely decided by the owner of the record'' \cite{Megyesi:2019}, and we met records marked partly decrypted that had been printed in full a century earlier. Computational work reports token accuracy and calls about 95--98\,\% solved without fixing a cut-off \cite{Ravi:2011}. Historians call a letter deciphered when the running text reads and a few names or code groups stay open. We state our own rule. A target is \emph{complete} when (1) at least 95\,\% of its cipher tokens give sense, counted from the reading file; (2) no remaining gap was left untried; (3) every document in it is covered, so that a group with one unread letter is partial; and (4) what stays open is scattered code groups or names, or pieces blocked from outside (no key material, too little text, illegible, or held only in the archive). Conditions (2) and (4) are ours. Open nomenclator codes are counted separately, since a letter can read at 99\,\% of tokens with many code entries unknown. Where the share of tokens with a value and the share that reads as sense differ, both are given.

The profiles were written by model agents instructed to source every value from the target's folder and to measure lengths with the common script; Section~\ref{sec:validity} reports how accurate they are.

\section{Working method}
\label{sec:method}

\paragraph{The record.} The agents kept a git repository with one folder per target. Of the \Ntargets\ targets, \Nnotes\ have a notes file that the agent maintained while working: what was fetched and from where, the transcription conventions, each hypothesis and solver run with its result, the controls, the literature searched, the reading, and what remains open with the reason. Solver code is in the folder for \Ncode\ targets; transcriptions and reading files sit beside it, with crops of the manuscript images where their terms allow it. A finished target was also written up for a public web page and entered in shared ledgers of results. The agents committed their work as they went, and \Ncommitsmodel\ of the \Ncommits\ commits name the Claude model that made them in a co-author line (the rest are merges, a few housekeeping commits by hand, and commits from the Codex harness, whose model is recorded in the profile instead), so the history fixes the order and date of the steps. The profiles of Section~\ref{sec:schema} were filled from these folders and this history.

\paragraph{Models and harness.} The work was done by LLMs running as coding agents in terminal harnesses with file, shell and web access. Claude models (mainly Opus~5, Opus~5.5 and Fable~5.1) worked on \Nclaude\ targets in Claude Code; OpenAI GPT-5 and GPT-6 models worked on \Ngpt\ targets in the Codex harness, \Nboth\ of them shared with a Claude model; the model is not recorded for \Nmodelunknown. Long tasks were split among parallel sub-agents, for example one per manuscript page for transcription, and merged by the main agent. The corpus is too unbalanced between model families to compare them, and we do not.

\paragraph{What the humans did.} The first author's part was small and mostly at the project level. The first author obtained a DECODE login with the privileges needed to see images and keys and handed the agent the session cookie, so that it could fetch records itself; set up the harness; and posed each problem as a standing goal, which keeps the agent working until the stated condition holds. During the work the first author occasionally chose which of several leads the agent had identified to pursue. The first author did not transcribe, propose keys or readings, or write code. Outside contributors sent in solutions as pull requests, issues and emails, which the first author reviewed against the images; those solutions are the contributors' and are reported apart (Section~\ref{sec:contrib}). On \Nexternal\ targets the work was shared between the project and an outside solver (mixed provenance); they stay in the figures and are marked $^\dagger$ in the tables and the appendix: \ContribTargets. For these the profile records the agent's steps, and the contributor's part is credited to them in the ledger. In \Nhumaninput\ targets the first author went further and supplied a specific source (a printed book the agent could not download) or suggested a route (a candidate key, archive addresses); these are marked in the data. The agents did the rest: locating images through archive interfaces, transcription, statistics, hypotheses, writing and running solvers, searching printed editions, and writing up.

\paragraph{Tools.} Solvers were written per target, mostly in Python with C\# for speed: hill climbing and simulated annealing over homophonic and syllabic keys, annealers constrained to an alphabetical key design, monotone alignment of ordered key slots, hard EM aligners from ciphertext to a contemporary decipherment, and SAT encodings. The scoring models were character $n$-gram models built from period corpora, for example nineteenth-century editions of sixteenth-century French correspondence, shared across targets through one registry.

\paragraph{Conventions.} Three project rules bear on the numbers. A negative result is reported only alongside a matched control: a synthetic ciphertext of the same length, alphabet and design that the same solver recovers, or fails to recover in the same way. Every reading is graded by its source: a primary key, a known-plaintext letter, uncertain, or inferred from context or from the alphabetical position of a sign. And before a target was treated as unsolved, the nineteenth-century editions of the correspondence and the DECODE records of the same volume were searched.

\section{The corpus}
\label{sec:corpus}

\begin{table}[t]
\centering\small
\begin{tabular}{@{}lrrr@{}}
\toprule
outcome method & $n$ & complete & partial \\
\midrule
key from ciphertext only & 43 & 16 & 21 \\
key from external plaintext & 12 & 6 & 6 \\
key from adjacent plaintext & 47 & 18 & 27 \\
\quad\emph{key recovered, all routes} & \emph{102} & \emph{40} & \emph{62} \\
external key matched & 26 & 8 & 13 \\
known key applied & 88 & 42 & 36 \\
existing decipherment & 104 & 89 & 14 \\
\midrule
not solved & 114 & -- & -- \\
\midrule
not applicable & 49 & -- & -- \\
\midrule
all targets & 483 & 179 & 117 \\
\bottomrule
\end{tabular}
\caption{\label{tab:methods} Outcome method and extent for the 483 targets. The first three methods recover the key here; the next two apply a key that already existed; the sixth transcribes and checks a decipherment that already existed. Complete means the read bar of Section~\ref{sec:schema} is met.}
\end{table}

\begin{table}[t]
\centering\small
\begin{tabular}{@{}lr@{\qquad}lr@{}}
\toprule
century & $n$ & plaintext language & $n$ \\
\midrule
14th c. & 3 & French & 120 \\
15th c. & 34 & English & 88 \\
16th c. & 183 & Italian & 80 \\
17th c. & 93 & German & 48 \\
18th c. & 94 & Spanish & 43 \\
19th c. & 28 & Latin & 34 \\
20th c. & 28 & Dutch & 14 \\
undated & 20 & Polish & 9 \\
\bottomrule
\end{tabular}
\caption{\label{tab:corpus} The corpus by century of the earliest document and by plaintext language (a target with letters in two languages counts under both; 17 languages in all).}
\end{table}

\begin{table}[t]
\centering\small
\resizebox{\columnwidth}{!}{%
\begin{tabular}{@{}lrrrrr@{}}
\toprule
region (most documents) & $n$ & recovered & applied & existed & not solved \\
\midrule
\multicolumn{6}{@{}l}{\emph{project targets}} \\
\quad Europe & 390 & 90 & 106 & 96 & 98 \\
\quad North America & 16 & 6 & 1 & 5 & 4 \\
\quad East Asia & 2 & 1 & 1 & 0 & 0 \\
\quad Middle East and South Asia & 3 & 1 & 0 & 2 & 0 \\
\quad not recorded & 23 & 4 & 6 & 1 & 12 \\
\quad\emph{all} & 434 & 102 & 114 & 104 & 114 \\
\midrule
\multicolumn{6}{@{}l}{\emph{contributed targets}} \\
\quad Europe & 113 & 65 & 31 & 14 & 3 \\
\quad North America & 14 & 3 & 7 & 4 & 0 \\
\quad Latin America & 3 & 2 & 1 & 0 & 0 \\
\quad East Asia & 18 & 16 & 1 & 1 & 0 \\
\quad Oceania & 1 & 1 & 0 & 0 & 0 \\
\quad not recorded & 23 & 20 & 3 & 0 & 0 \\
\quad\emph{all} & 172 & 107 & 43 & 19 & 3 \\
\bottomrule
\end{tabular}}
\caption{\label{tab:region} Outcome by region of the correspondence: the region named by most documents of the target (Europe where no other is named). Upper block: the project's targets; lower block: the contributed targets of Section~\ref{sec:contrib}. Columns as in Table~\ref{tab:attack}.}
\end{table}

\begin{table}[t]
\centering\small
\resizebox{\columnwidth}{!}{%
\begin{tabular}{@{}lrrrrr@{}}
\toprule
starting point & $n$ & recovered & applied & existed & not solved \\
\midrule
\multicolumn{6}{@{}l}{\emph{project targets}} \\
\quad ciphertext only & 110 & 38 & 4 & 4 & 64 \\
\quad crib from clear text & 30 & 22 & 1 & 3 & 4 \\
\quad decipherment or known plaintext & 83 & 26 & 6 & 43 & 8 \\
\quad a key (archive, print, sibling) & 139 & 14 & 102 & 14 & 9 \\
\quad none recorded & 72 & 2 & 1 & 40 & 29 \\
\quad\emph{all} & 434 & 102 & 114 & 104 & 114 \\
\midrule
\multicolumn{6}{@{}l}{\emph{contributed targets}} \\
\quad ciphertext only & 52 & 51 & 0 & 0 & 1 \\
\quad crib from clear text & 4 & 4 & 0 & 0 & 0 \\
\quad decipherment or known plaintext & 62 & 49 & 0 & 13 & 0 \\
\quad a key (archive, print, sibling) & 51 & 3 & 42 & 6 & 0 \\
\quad none recorded & 3 & 0 & 1 & 0 & 2 \\
\quad\emph{all} & 172 & 107 & 43 & 19 & 3 \\
\bottomrule
\end{tabular}}
\caption{\label{tab:attack} Outcome by what the attack started from (profile field conditions.attack). Upper block: the project's targets; lower block: the contributed targets of Section~\ref{sec:contrib}. Columns: key recovered (from ciphertext only, adjacent or external plaintext); key applied (a known key, or an external key matched to the document); a decipherment already existed; not solved. Targets classed not applicable are left out.}
\end{table}

\begin{table}[t]
\centering\small
\resizebox{\columnwidth}{!}{%
\begin{tabular}{@{}lrrrrr@{}}
\toprule
cipher family & $n$ & recovered & applied & existed & not solved \\
\midrule
\multicolumn{6}{@{}l}{\emph{project targets}} \\
\quad homophonic or simple + nomenclator & 179 & 45 & 66 & 47 & 21 \\
\quad code or codebook & 72 & 15 & 9 & 16 & 32 \\
\quad undetermined & 66 & 2 & 3 & 16 & 45 \\
\quad syllabic nomenclator & 49 & 10 & 20 & 12 & 7 \\
\quad substitution, no nomenclator & 48 & 24 & 11 & 9 & 4 \\
\quad transposition, polygraphic, polyalphabetic & 20 & 6 & 5 & 4 & 5 \\
\quad\emph{all} & 434 & 102 & 114 & 104 & 114 \\
\midrule
\multicolumn{6}{@{}l}{\emph{contributed targets}} \\
\quad homophonic or simple + nomenclator & 65 & 42 & 16 & 7 & 0 \\
\quad code or codebook & 23 & 4 & 11 & 7 & 1 \\
\quad undetermined & 28 & 19 & 5 & 3 & 1 \\
\quad syllabic nomenclator & 21 & 13 & 8 & 0 & 0 \\
\quad substitution, no nomenclator & 14 & 13 & 1 & 0 & 0 \\
\quad transposition, polygraphic, polyalphabetic & 21 & 16 & 2 & 2 & 1 \\
\quad\emph{all} & 172 & 107 & 43 & 19 & 3 \\
\bottomrule
\end{tabular}}
\caption{\label{tab:family} Outcome by cipher family, from the system types in the profile. Upper block: the project's targets; lower block: the contributed targets of Section~\ref{sec:contrib}. Columns as in Table~\ref{tab:attack}.}
\end{table}

\begin{table}[t]
\centering\small
\resizebox{\columnwidth}{!}{%
\begin{tabular}{@{}lrrrrr@{}}
\toprule
prior reading & $n$ & recovered & applied & existed & not solved \\
\midrule
\multicolumn{6}{@{}l}{\emph{project targets}} \\
\quad none known & 129 & 41 & 23 & 0 & 65 \\
\quad existed, known before the attempt & 151 & 24 & 51 & 62 & 14 \\
\quad existed, found during the attempt & 100 & 21 & 30 & 39 & 10 \\
\quad existed, found after or not at all & 20 & 11 & 6 & 1 & 2 \\
\quad existed, when found unknown & 3 & 0 & 1 & 2 & 0 \\
\quad unknown & 31 & 5 & 3 & 0 & 23 \\
\quad\emph{all} & 434 & 102 & 114 & 104 & 114 \\
\midrule
\multicolumn{6}{@{}l}{\emph{contributed targets}} \\
\quad none known & 81 & 59 & 21 & 0 & 1 \\
\quad existed, known before the attempt & 23 & 9 & 10 & 4 & 0 \\
\quad existed, found during the attempt & 19 & 3 & 3 & 11 & 2 \\
\quad existed, found after or not at all & 46 & 33 & 9 & 4 & 0 \\
\quad unknown & 3 & 3 & 0 & 0 & 0 \\
\quad\emph{all} & 172 & 107 & 43 & 19 & 3 \\
\bottomrule
\end{tabular}}
\caption{\label{tab:prior} Outcome by whether a reading of the text existed anywhere before this project, and when it was found (the contamination question). Upper block: the project's targets; lower block: the contributed targets of Section~\ref{sec:contrib}. Columns as in Table~\ref{tab:attack}.}
\end{table}

\begin{table}[t]
\centering\small
\resizebox{\columnwidth}{!}{%
\begin{tabular}{@{}lrrrrr@{}}
\toprule
ciphertext tokens & $n$ & recovered & applied & existed & not solved \\
\midrule
\multicolumn{6}{@{}l}{\emph{project targets}} \\
\quad $<$ 250 & 113 & 23 & 22 & 25 & 43 \\
\quad 250--999 & 119 & 27 & 34 & 19 & 39 \\
\quad 1,000--4,999 & 89 & 33 & 29 & 16 & 11 \\
\quad $\geq$ 5,000 & 31 & 12 & 14 & 4 & 1 \\
\quad unknown & 82 & 7 & 15 & 40 & 20 \\
\quad\emph{all} & 434 & 102 & 114 & 104 & 114 \\
\midrule
\multicolumn{6}{@{}l}{\emph{contributed targets}} \\
\quad $<$ 250 & 44 & 24 & 11 & 8 & 1 \\
\quad 250--999 & 71 & 47 & 18 & 5 & 1 \\
\quad 1,000--4,999 & 45 & 33 & 9 & 3 & 0 \\
\quad $\geq$ 5,000 & 6 & 2 & 4 & 0 & 0 \\
\quad unknown & 6 & 1 & 1 & 3 & 1 \\
\quad\emph{all} & 172 & 107 & 43 & 19 & 3 \\
\bottomrule
\end{tabular}}
\caption{\label{tab:length} Outcome by ciphertext length (enciphered tokens summed over the documents of a target). Upper block: the project's targets; lower block: the contributed targets of Section~\ref{sec:contrib}. Columns as in Table~\ref{tab:attack}.}
\end{table}

\begin{table*}[t]
\centering\small
\resizebox{\textwidth}{!}{%
\begin{tabular}{@{}lrllrll@{}}
\toprule
target & year & language & system & tokens & key order & extent \\
\midrule
siena14 & 1420 & Italian & homophonic or simple + nomenclator & 1291 & random & none \\
r1854 & 1428 & Latin & homophonic or simple + nomenclator & 1802 & not applicable & complete \\
cattanei1430 & 1430 & Latin & substitution, no nomenclator & 1736 & not applicable & partial \\
siena25 & 1450 & Italian & substitution, no nomenclator & 496 & random & complete \\
siena18$^\dagger$ & 1496 & Latin & homophonic or simple + nomenclator & 1469 & random & complete \\
esp318no95$^\dagger$ & 1497 & Spanish & homophonic or simple + nomenclator & 1364 & unknown & partial \\
vich1515 & 1515 & Spanish & homophonic or simple + nomenclator & 3177 & unknown & none \\
soria1523 & 1523 & Spanish & homophonic or simple + nomenclator & -- & unknown & partial \\
ranzo1528$^\dagger$ & 1528 & Italian & code or codebook & 8224 & random & partial \\
rangone1530 & 1529 & Italian & substitution, no nomenclator & 264 & unknown & partial \\
joachim1530 & 1530 & Italian & homophonic or simple + nomenclator & 402 & unknown & partial \\
guiche1551 & 1551 & French & substitution, no nomenclator & 431 & unknown & partial \\
memoire1560 & 1560 & French & code or codebook & 1715 & unknown & partial \\
r9424 & 1583 & German & homophonic or simple + nomenclator & 913 & random & partial \\
harley1582r8505 & 1585 & French & substitution, no nomenclator & 3564 & random & none \\
segur & 1585 & French & syllabic nomenclator & 466 & alphabetical & complete \\
lorraine1592$^\dagger$ & 1592 & French & homophonic or simple + nomenclator & 1109 & unknown & none \\
stowe166 & 1593 & English & homophonic or simple + nomenclator & 126 & unknown & complete \\
segura1596 & 1596 & Spanish & substitution, no nomenclator & 125 & unknown & complete \\
senisteros1605 & 1605 & Spanish & substitution, no nomenclator & 1447 & unknown & complete \\
anhalt1610 & 1610 & German & undetermined & 242 & not applicable & none \\
ferdinand1619 & 1619 & German & homophonic or simple + nomenclator & 4181 & partly alphabetical & complete \\
warsaw & 1627 & Italian & syllabic nomenclator & 465 & alphabetical & none \\
harley286 & 1635 & Latin & substitution, no nomenclator & 281 & random & partial \\
marburg1635 & 1635 & German & substitution, no nomenclator & 590 & unknown & partial \\
malsburg1637$^\dagger$ & 1637 & French, German & transposition, polygraphic, polyalphabetic & 13487 & not applicable & partial \\
transylvania1645 & 1645 & Hungarian & undetermined & 4785 & unknown & partial \\
mercy1648$^\dagger$ & 1648 & Spanish & homophonic or simple + nomenclator & 529 & random & partial \\
wesselenyi1664 & 1664 & Latin & homophonic or simple + nomenclator & 244 & alphabetical & complete \\
schenck1706 & 1706 & German & homophonic or simple + nomenclator & 1181 & alphabetical & partial \\
potocka1714 & 1715 & French, Polish & homophonic or simple + nomenclator & 886 & partly alphabetical & partial \\
kauderbach1754 & 1754 & French & homophonic or simple + nomenclator & 8512 & random & partial \\
swieten1759 & 1759 & French & homophonic or simple + nomenclator & 553 & alphabetical & complete \\
r1892 & 1795 & Dutch, French & substitution, no nomenclator & 2067 & random & partial \\
mountflorence1801 & 1801 & English & code or codebook & 39 & alphabetical & complete \\
soglia1848 & 1848 & Italian & code or codebook & 955 & alphabetical & complete \\
augur1874 & 1874 & English & transposition, polygraphic, polyalphabetic & 246 & not applicable & partial \\
sunyatsen & 1916 & Chinese & transposition, polygraphic, polyalphabetic & 177 & partly alphabetical & complete \\
lejava1921 & 1921 & Ru & transposition, polygraphic, polyalphabetic & 167 & unknown & complete \\
voikov1921 & 1921 & Ru & transposition, polygraphic, polyalphabetic & 615 & unknown & complete \\
egmond & -- & French & substitution, no nomenclator & 774 & unknown & complete \\
r9416 & -- & German & substitution, no nomenclator & 3313 & unknown & partial \\
sp106r927 & -- & French & substitution, no nomenclator & 1348 & unknown & partial \\
\bottomrule
\end{tabular}}
\caption{\label{tab:ct} The 43 targets whose key was recovered from the ciphertext alone. $^\dagger$~solved in part by an outside contributor.}
\end{table*}

\begin{table}[t]
\centering\small
\begin{tabular}{@{}lrrr@{}}
\toprule
outcome method & $n$ & complete & partial \\
\midrule
key from ciphertext only & 52 & 28 & 22 \\
key from external plaintext & 9 & 6 & 3 \\
key from adjacent plaintext & 46 & 17 & 28 \\
external key matched & 9 & 2 & 6 \\
known key applied & 34 & 15 & 18 \\
existing decipherment & 19 & 9 & 10 \\
\midrule
not solved & 3 & -- & -- \\
\midrule
not applicable & 1 & -- & -- \\
\midrule
all contributed targets & 173 & 77 & 87 \\
\bottomrule
\end{tabular}
\caption{\label{tab:contrib} Outcome method and extent for the 173 contributed targets: those whose solution is credited to someone outside the project. The model figures of Table~\ref{tab:methods} leave them out.}
\end{table}


The corpus is restricted to historical (classical) ciphers dated between 1350 and 1950, from any region; the project's work on later or modern material is not reported here. Table~\ref{tab:corpus} summarises the \Ntargets\ targets. They hold \Ndocs\ ciphertext documents, from \YearMin\ to \YearMax. Most are diplomatic letters of the sixteenth (\NsixteenthC) and seventeenth (\NseventeenthC) centuries: French, Spanish, Italian, English, Imperial and papal embassy traffic, with \NfifteenthC\ targets from the fifteenth century and \NmodernC\ from the nineteenth and twentieth, including telegrams and military radiograms. Another \Nundated\ targets have no established date but come from early-modern archival contexts. The main plaintext languages are \TopLanguages.

A length is known for \Nlengthknown\ targets (measured by script on every document for \Nmeasured); together they hold \TokensTotal\ enciphered tokens. The agent transcribed \Ndocsllm\ of the documents (\PctDocsllm\,\%) itself from archive images; \Ndocsdecode\ came from DECODE transcriptions and \Ndocspublished\ from published ones.

In all, \NNA\ targets are classed not applicable: \NnotCipher\ turned out not to be ciphers (clear text with a cipher instruction, shorthand, key sheets, dummy numbers) and \NoffLine\ could not be obtained because the manuscript was not digitised or not reachable. The analysis below uses the \Nworked\ targets that were worked.

\section{Results}
\label{sec:results}

\paragraph{Overall.} Table~\ref{tab:methods} gives the outcome for every target. The text was obtained for \Nobtained\ of the \Nworked\ worked targets: \Ncomplete\ complete and \Npartial\ partial; the remaining \NNS\ were not solved. The methods behind the \Nobtained\ differ widely. The key was recovered here for \Nrecovered\ targets (\Nrecoveredcomplete\ complete): \NCT\ from the ciphertext alone, \NAPT\ from adjacent plaintext and \NEPT\ from external plaintext. For \Napplied\ targets an existing key was applied: \NKK\ with a key already tied to the correspondence and \NEK\ with a key found elsewhere and matched to the document. For \Ndeciph\ targets a complete decipherment already existed, and the contribution was a transcription, a check of the decipherment, or a correction of the catalogue. Counting only the first three methods as cryptanalysis, the model recovered \Nrecovered\ keys in \Nworked\ attempts.

\paragraph{Starting point dominates.} Table~\ref{tab:attack} sorts the worked targets by what the attack started from. With a key in hand, from an archive, a printed edition or a sibling letter of the same office, the text was obtained in \Nkeyedobtained\ of \Nkeyed\ cases (\Nkeyedpct\,\%). With a contemporary decipherment of some of the letters, it was obtained in \Ndecstartobtained\ of \Ndecstart; most of these rebuilt the key from the decipherment and applied it to the undeciphered letters. With a crib from the clear text around the cipher, it was obtained in \Ncribobtained\ of \Ncrib. From the ciphertext alone, the model recovered the key in \NctoCT\ of \Ncto\ attempts (\NctoPct\,\%); in \NctoDEC\ more, a decipherment turned up during the attempt, and \NctoNot\ were not solved. The other \NCTnotcto\ ciphertext-only recoveries began from a crib guessed from the subject of the letter or from the plaintext vocabulary of a sibling key, which uses no plaintext of the document itself.

\paragraph{The ciphertext-only recoveries.} Table~\ref{tab:ct} lists the \NCT\ targets whose key was recovered from the ciphertext alone. Most are homophonic or simple substitutions, some with a small nomenclator, on a few hundred to a few thousand tokens. The key order is known for \NCTalpha\ of them to be alphabetical or partly alphabetical and for \NCTrandom\ to be random; an alphabetical key let the model predict unseen values and confirm them in the text, and in one case (Warsaw 1627) the key's alphabetical order, which the solver did not reward, was the independent check. All but \the\numexpr\NCT-\NCTverified\relax\ of the \NCT\ carry a verification beyond the reading itself: a matched control, key consistency, or historical consistency with the known events. Of the \NCT, \NCTprior\ had a reading somewhere that was found after the key was rebuilt or that covered only part of the target (\NCTpriorbefore\ known before the attempt, covering other letters or passages than those whose key was rebuilt); they are independent re-solutions and are marked as such in the data.

\paragraph{Prior readings.} Table~\ref{tab:prior} splits the worked targets by whether a reading existed anywhere before the project, and when it was found. Where none is known (\Nclean\ targets), the text was obtained for \Ncleanobtained, \Ncleancomplete\ of them complete: \Ncleanrecovered\ by recovering the key (\NcleanCT\ from the ciphertext alone) and \NcleanApplied\ by applying a key that had never been used on these letters. This is the subset free of contamination, since the model cannot have memorised a decipherment that does not exist. Where a reading existed, it was most often found by the agent during the attempt (\NpriorFoundDuring\ targets) rather than known beforehand (\NpriorFoundBefore); \Nprintprior\ targets had a reading in print, usually in a nineteenth-century edition that the modern lists of unsolved ciphers had missed.

\paragraph{Cipher family matters less.} By family (Table~\ref{tab:family}), homophonic ciphers with nomenclators are the largest group and were obtained most often, but most of those attacks started from a key. Simple and homophonic substitutions without a nomenclator have the most key recoveries relative to their number. The \emph{undetermined} row, targets whose system was never established, is mostly not solved by construction. The comparison across families is confounded by starting point, and the cells are small.

\paragraph{Length.} Short ciphertexts (under 250 tokens) were most often not solved (Table~\ref{tab:length}); above that, the share not solved falls with length, and almost no text of 5,000 tokens or more went unsolved. Long texts are mostly nomenclator letters for which a key was found or rebuilt from sibling decipherments.

\paragraph{Effort.} The profiles hold \Nsteps\ solution steps, of which \Nfailedsteps\ (\PctFailedSteps\,\%) failed or were ruled out. There are \Nsolversteps\ solver runs, of which \Nsolverfailed\ failed or were ruled out, \Nlitsteps\ literature searches and \Ntranssteps\ transcription steps. Matched controls appear in \Ncontrol\ targets. The balance of steps reflects the work: searching the literature and the archives was as frequent as running solvers, and more often decisive.

\paragraph{What else the work added.} Beyond the text, the profiles record corrections to the catalogues (\Ncatcorr\ targets: a wrong sender, date, or decryption status in DECODE or on a list of unsolved ciphers), identifications of an unknown sender or date (\Nsenderid), keys extended or corrected (\Nkeyextended), and \Nambiguous\ decipherments with a polyphonic or variable-length key where applying a known key required judgement.

\section{Contributed solutions}
\label{sec:contrib}

The project site invites solutions, new ciphers and checks from anyone, credits each by name in a ledger, and keeps contributed results out of the model figures above. Between \ContribDateFirst\ and \DateLast, \Ncontributors\ contributors, with the permission recorded for this paper, added \Ncontrib\ targets to the ledger (Table~\ref{tab:contrib}), dated \NcontribYearMin--\NcontribYearMax\ and written in \NcontribLanguages\ languages. The text was obtained for \Ncontribobtained\ of the \Ncontribworked\ that could be worked (\Ncontribcomplete\ complete, \Ncontribpartial\ partial). The key was recovered by the contributor for \Ncontribrecovered: \NcCT\ from the ciphertext alone, \NcAPT\ from adjacent plaintext and \NcEPT\ from external plaintext; \NcEK\ and \NcKK\ applied an external or a known key, and \NcDEC\ transcribed an existing decipherment.

We do not fold these into the model figures, for two reasons. The solver is not the project's agent: the starting conditions, the prompts and the tools are the contributor's, and the profile can record them only as the contributor reported them. And the contributions are not uniform. The second author's \NcontribTop\ targets were solved with CrypTool's transcriber and solver, in some cases revised or read afterwards with a model; for most of the others the profile records the contributor directing a language model (Claude Code, Codex or the contributor's own agents), and for a few it records neither. The categories of Section~\ref{sec:schema} apply to a human solution as well as to a model's, since a key recovered from ciphertext only is that whoever recovered it, so the contributed targets carry the same profile and method fields. Table~\ref{tab:contributors} gives the targets by credited solver and Table~\ref{tab:allcontrib} in the appendix lists them. Figures in Section~\ref{sec:results} that cite ``the model'' refer to the \Ntargets\ project targets only.

\ifanon\else
The contributors are \ContribCredits.
\fi

\begin{table}[t]
\centering\small
\resizebox{\columnwidth}{!}{%
\begin{tabular}{@{}lrrrrrrrr@{}}
\toprule
contributor & $n$ & CT & APT & EPT & EK & KK & DEC & NS \\
\midrule
\Who{C1}{George Lasry} & 77 & 30 & 33 & 0 & 5 & 5 & 4 & 0 \\
\Who{C2}{Brice Dudout} & 50 & 20 & 4 & 6 & 3 & 14 & 2 & 1 \\
\Who{C3}{WAINUT} & 30 & 3 & 1 & 1 & 3 & 8 & 10 & 4 \\
\Who{C4}{Michael Brackx} & 10 & 2 & 2 & 0 & 0 & 3 & 3 & 0 \\
\Who{C5}{Arya Sanketbhai Patel} & 8 & 1 & 0 & 0 & 0 & 2 & 4 & 0 \\
\Who{C6}{Harumi-Matsuba} & 7 & 0 & 1 & 0 & 0 & 6 & 0 & 0 \\
\Who{C7}{Feyseel Nur} & 6 & 0 & 2 & 1 & 1 & 1 & 1 & 0 \\
\Who{C8}{alejandrozarco} & 2 & 0 & 1 & 0 & 0 & 1 & 0 & 0 \\
\Who{C9}{Angus Robinson} & 2 & 1 & 1 & 0 & 0 & 0 & 0 & 0 \\
\Who{C10}{Brett Burford} & 2 & 0 & 0 & 0 & 0 & 2 & 0 & 0 \\
\Who{C11}{Larry Beck} & 2 & 1 & 0 & 0 & 1 & 0 & 0 & 0 \\
\Who{C12}{Harrison Bornstein} & 1 & 0 & 0 & 1 & 0 & 0 & 0 & 0 \\
\Who{C13}{Paolo Rosson} & 1 & 0 & 0 & 0 & 0 & 1 & 0 & 0 \\
\Who{C14}{johonn} & 1 & 0 & 1 & 0 & 0 & 0 & 0 & 0 \\
\Who{C15}{Satoru} & 1 & 0 & 0 & 0 & 1 & 0 & 0 & 0 \\
\Who{C16}{setsunaatto} & 1 & 0 & 0 & 0 & 1 & 0 & 0 & 0 \\
\Who{C17}{NoAutopilot} & 1 & 1 & 0 & 0 & 0 & 0 & 0 & 0 \\
\bottomrule
\end{tabular}}
\caption{\label{tab:contributors} Targets by the contributor credited with the solution, for the contributed targets and the mixed ones of Table~\ref{tab:methods}. A target with two solvers counts under each. Method abbreviations as in Table~\ref{tab:all}.}
\end{table}


\section{Case studies}
\label{sec:cases}

\paragraph{Toledo to Philip II, 1565 (key from adjacent plaintext).} A letter of the viceroy of Sicily during the siege of Malta, with 1,937 two-digit figures inline in clear Spanish (Archivo General de Simancas, Estado leg.~1394). The agent worked out the image route of the Spanish archives portal, transcribed the figures, and took two runs next to clear words as cribs (\emph{seiscientos}, \emph{en tiera}). They fixed eleven letters, and the values came out in alphabetical order, so the remaining letters were predicted from their neighbours and confirmed in words. Nonsense in the decode located misread digits, which were corrected on the images. A companion letter of the same day, found by searching the portal, decodes with the same key. No earlier decipherment is known.

\paragraph{Balbases to Fuenmayor, 1677--78 (existing decipherment; key rebuilt from adjacent plaintext).} Fourteen Spanish despatches from the peace congress at Nijmegen, ten with a contemporary decipherment in the margin. Eight keys of the same series were tried and none fitted. Two sub-agents transcribed every cipher and margin line; anchors from $n$-grams repeated across letters seeded a hard-EM aligner that maps each token to a substring of the margin text, and the table converged in two rounds. It reproduces the ten margins. Applied to the four letters that carry no margin, the table gives a syllable stream with a few groups still without a value. A reading audit of 22 October counted those four as decoded, not read, so only the ten margined letters, 71\,\% of the tokens of the group, count as read, and the profile classes the target by the margins and as partial.

\paragraph{Feuqui\`eres to Catinat, 1691 (key from external plaintext).} A two-part code of 418 groups (160 distinct) in the ``petit chiffre'' of the governors of Pignerol, deciphered by \'Etienne Bazeries in 1893 without publishing the table. A matched control of the inferred design showed that annealing could not recover it: wrong keys outscored the true one. The agent then searched the digitised text of the 1819 \emph{M\'emoires} of Catinat for runs of numbers and found a letter of Louvois of 6~September 1690 printed in the same code with its translation. Anchoring by hand and propagating between the two letters gave a value to 586 of 601 tokens, and the Feuqui\`eres text agrees with Catinat's own memoir of the same days. The reading audit counts only the Louvois letter (183 tokens) as read, since the Feuqui\`eres letter, with twelve groups still open, is decoded but its sense is only half established; the target is partial (30\,\% of tokens read).

\paragraph{Bordeaux to Brienne, 1653 (external key matched).} A French despatch intercepted in London: 813 tokens of graphic signs and numbers, three number series marked by a prime, two dots or an overbar. From sibling keys of the same office the agent inferred the design (alphabetical syllabaries in the marked series) and built an annealer with that constraint. It recovered 99\,\% of matched controls and failed on the real text; controls with 10\,\% of the diacritics misread failed in the same way (52\,\% and 9\,\%). The agent then found the English Deciphering Branch key sheet for this correspondent in a separate volume of the British Library (DECODE R7537). It confirmed the design exactly, and a check against the images made 21 corrections to the transcription, splitting two merged glyphs and restoring lost marks. The key applied to the corrected transcription read the letter except name codes and rare tokens.

\paragraph{Henry of Navarre to S\'egur, 1585--86 (key from ciphertext only).} Figures in three letters, 466 tokens. Before any value was known, the frequencies of the upper figures taken modulo five fell into blocks of five, the signature of an alphabetical syllabary of consonants by vowels. An annealer constrained to that design, with the clear French around the cipher as fixed context, converged on a key and recovered 96.6\,\% of matched controls under random keys of the same shape. A sibling key published by Tomokiyo, with different offsets, read a fourth letter; the constrained annealer then re-derived that fourth key blind at 93\,\% of letter tokens, against 90\,\% on a matched control.

\paragraph{Ordnungspolizei radiograms, 1942 (not solved).} Seven parts of 27~February 1942 in the German police Doppelkasten, a double Playfair in which each pair passes twice through the same two boxes, 970 letters under one key. The agent implemented the system from the German manuals and attacked it with annealing, iterated local search, tempering, relaxation and SAT. All of these failed on synthetic texts of the same length with a known key: the double pass leaves a local search no gradient. A known-plaintext solver it wrote recovers the boxes from 80 to 111 plaintext pairs on synthetic data, but no plaintext of these messages is available.

\section{Where the models failed}
\label{sec:failures}

\paragraph{Transcription.} The largest single cause of failure was the transcription, not the cryptanalysis. The agent transcribed every document of \Nimg\ worked targets from images and obtained the text for \Nimgobtained\ of them, but the Bordeaux case shows how a transcription that looks usable can defeat a correct solver. Controls with injected transcription noise are cheap and diagnostic, and we recommend them as a standard step.

\paragraph{Short texts and large nomenclators.} Large nomenclators on a few hundred tokens were beyond reach without a key or a crib, as they are for human analysts. The agent's value here was in finding the key: on the DECODE record of a sibling letter, in the same archive volume, or in a printed edition.

\paragraph{Undetermined systems.} Where the system was never identified, the target was almost never solved. These are mostly short texts, from a few lines of figures to a few hundred tokens: an isolated slip, a postscript, a fragment of a series whose other parts are lost.

\paragraph{Double encipherment and hand systems of the twentieth century.} The double Playfair, Enigma traffic and most of the other twentieth-century hand systems and codes in the corpus were not solved, while three short transposition and code targets of 1916--1921 had their keys recovered from the ciphertext alone. The agents wrote correct implementations and competent searches for the failures and reported the negative results with controls.

\paragraph{Finding what was already known.} Many targets on current lists of unsolved ciphers had been read long ago, most often in a nineteenth-century printed edition, sometimes in a decipherment bound in the next folio. The agent found these by searching the editions and the neighbouring archive records, and we classify them by the method that applied, never as key recoveries.

\section{Threats to validity}
\label{sec:validity}

\paragraph{The profiles are written by the agents.} They were compiled by model agents from the project's notes, partly during the work and partly afterwards, and the key order is unknown for \Nkeyorderunknown\ targets, usually because the key was never recovered. The human's role is recorded at the project level for most targets, since the notes rarely mention it; lead choices within a session are not itemised. To estimate accuracy, ten profiles were drawn at random (fixed seed) from the first 82 and checked field by field against the notes, the result tables and the ciphertext files by fresh agents that had not written them. Of 249 fields checked, 38 were unknown. Of the 211 with a value, 196 agreed with the sources (93\,\%), 5 contradicted them and 10 asserted something the sources do not establish. All recorded lengths re-measured exactly. The errors were of three kinds: placeholder values in profiles of targets never obtained; a modern partial reading filed as a plaintext location; and over-confident system fields for ciphers that were never broken. They were corrected, and the schema gained a \emph{margin} value for plaintext location. The audit was done by agents, not by a human, which limits what it shows.

\paragraph{The outcome method.} The method was added to every profile after the work, by review agents that read each target's notes. A rule-based proposal computed from the older profile fields (attack class, prior solution, solution steps) agreed with the reviewers' method for \Nreviewagree\ of \Nreview\ targets (\PctReviewagree\,\%). The proposal flagged \Nreviewflagged\ targets as ambiguous, typically whether a sibling key had already been tied to the document (key applied) or had to be matched to it (external key), and whether a key came from plaintext beside the cipher or from plaintext elsewhere; \Nflaggedchanged\ of the \Nreviewchanged\ disagreements fell among those. The boundary between \emph{known key applied} and \emph{external key matched} is the least certain; the boundary between recovering a key and applying one is well supported by the notes.

\paragraph{The extent rule is ours.} The 95\,\% threshold follows computational practice, not a published standard. A numeric share of tokens read is recorded for \Nfrac\ of the \Nobtained\ targets with text (\Nfracmeasured\ measured from the reading files, the rest estimated from the notes). Raising the threshold to 98\,\% would move at most \Nbarup\ complete targets to partial; lowering it to 90\,\% would move at most \Nbardown\ partial targets to complete. Another \Npartialhigh\ targets read above 95\,\% but are partial for another reason, usually a document of the group left unread.

\paragraph{Selection.} The targets came from lists of unsolved ciphers, from DECODE records marked non-decrypted, from sweeps of library and archive catalogues and portals, and, for the contributed targets, from the contributors' own choice, and the agent chose the order of attack. Targets that looked hopeless were attempted briefly or set aside, so the not-solved count mixes serious failures and brief closures.

\paragraph{No fixed budget.} Sessions varied in length, model and operator, and the first author supplied logins and in a few cases a source. The record is of what happened, not of a controlled experiment.

\paragraph{Contamination.} A model may have seen a printed decipherment in training. The prior-reading field addresses this directly, and the \Nclean\ targets with no known prior reading are the clean subset. For the others, the profile records whether the work used the prior reading, and the method records whether the key was rebuilt or applied.

\section{Conclusion}

On historical ciphers, an agentic LLM behaves less like a solver than like a research assistant that can also write solvers. It reads the manuscript, finds the key or the decipherment in the archive or in print, recognises a crib in the clear text, and exploits structure such as an alphabetical key. Of the \Nobtained\ targets whose text it obtained, the key was recovered here for \Nrecovered, and from the ciphertext alone for \NCT. Its failures follow the classical limits: short texts, large nomenclators, double encipherment and noisy transcriptions. Classifying outcomes by method keeps these two contributions apart, and we recommend it to anyone reporting decipherments made with such tools. The profile schema, the \Ntargets\ profiles and the working record behind them (notes, transcriptions, code and commit history) are released with the paper so that later work can check any profile against its sources and add targets and models on the same terms.

\ifanon\else
\section*{Acknowledgements}
The contributed targets of Section~\ref{sec:contrib} and the shared targets marked $^\dagger$ are the work of the contributors named there, who agreed to their inclusion in this paper: \ContribCredits\ (targets credited). We thank them.
\fi

\bibliographystyle{histocrypt}
\bibliography{refs}

\appendix
\onecolumn
{\scriptsize
\begin{longtable}{@{}llrlllr@{}}
\caption{\label{tab:all} Every target in the corpus, by date of the earliest document. Methods: CT key from ciphertext only; EPT key from external plaintext; APT key from adjacent plaintext; EK external key matched; KK known key applied; DEC existing decipherment; NS not solved; NA not applicable. $^\dagger$~solved in part by an outside contributor. Full profiles are in the data release.}\\
\toprule
target & year & language & system & method & extent & tokens \\
\midrule
\endfirsthead
\toprule
target & year & language & system & method & extent & tokens \\
\midrule
\endhead
\bottomrule
\endlastfoot
vico1398 & 1398 & -- & undetermined & NS & none & -- \\
benedict1399 & 1399 & Catalan, Latin & homophonic or simple + nomenclator & DEC & complete & -- \\
santacruz1399 & 1399 & Latin & undetermined & DEC & complete & -- \\
jafet1400 & 1400 & Latin & undetermined & DEC & complete & -- \\
exea1408 & 1408 & Latin & undetermined & DEC & complete & -- \\
florence1414 & 1414 & -- & substitution, no nomenclator & NA & n/a & -- \\
comuel1415 & 1415 & Latin & undetermined & DEC & complete & -- \\
comuel1416 & 1416 & Latin & undetermined & DEC & complete & -- \\
siena14 & 1420 & Italian & homophonic or simple + nomenclator & CT & none & 1291 \\
florence1429 & 1425 & Italian & homophonic or simple + nomenclator & KK & none & 99 \\
medici1425 & 1425 & Italian & undetermined & NA & n/a & -- \\
heidelberg597 & 1426 & German & substitution, no nomenclator & DEC & complete & -- \\
malatesta1428 & 1428 & Italian, Latin & undetermined & NS & none & 262 \\
mutiliana1428 & 1428 & Italian & homophonic or simple + nomenclator & APT & complete & 894 \\
r1854 & 1428 & Latin & homophonic or simple + nomenclator & CT & complete & 1802 \\
cattanei1430 & 1430 & Latin & substitution, no nomenclator & CT & partial & 1736 \\
maria1435 & 1435 & Catalan & undetermined & NS & none & 42 \\
siena8 & 1440 & Italian & homophonic or simple + nomenclator & DEC & complete & 166 \\
visconti1446 & 1446 & Italian & syllabic nomenclator & EK & partial & 16588 \\
barbavara1447 & 1447 & Italian & syllabic nomenclator & DEC & complete & -- \\
visconti1447 & 1447 & Italian & homophonic or simple + nomenclator & EK & complete & 722 \\
birago1448 & 1448 & Italian & homophonic or simple + nomenclator & DEC & complete & 513 \\
siena25 & 1450 & Italian & substitution, no nomenclator & CT & complete & 496 \\
borghesi1456 & 1456 & Italian & homophonic or simple + nomenclator & APT & none & 144 \\
knorre1462 & 1462 & German & undetermined & NA & n/a & -- \\
beatrice1482 & 1482 & Italian & undetermined & DEC & complete & -- \\
matthias1482 & 1482 & Latin & homophonic or simple + nomenclator & DEC & complete & 330 \\
sadoleto1482 & 1482 & Italian & homophonic or simple + nomenclator & APT & partial & -- \\
beatrice1486 & 1486 & Italian & undetermined & DEC & complete & -- \\
beatrice1486b & 1486 & Italian & undetermined & DEC & complete & -- \\
buda1489 & 1489 & Italian & syllabic nomenclator & KK & partial & 469 \\
hm64 & 1490 & English, Latin & undetermined & NS & none & -- \\
costabili1493 & 1493 & Italian & not a cipher & NA & n/a & -- \\
siena18$^\dagger$ & 1496 & Latin & homophonic or simple + nomenclator & CT & complete & 1469 \\
esp318no95$^\dagger$ & 1497 & Spanish & homophonic or simple + nomenclator & CT & partial & 1364 \\
federico1497 & 1497 & Italian & homophonic or simple + nomenclator & DEC & complete & -- \\
garcilaso1498 & 1498 & Spanish & homophonic or simple + nomenclator & NS & none & 318 \\
esp318 & 1500 & Spanish & homophonic or simple + nomenclator & KK & none & -- \\
ptr4715 & 1500 & -- & undetermined & NA & n/a & -- \\
estrada1503 & 1503 & Spanish & homophonic or simple + nomenclator & KK & partial & -- \\
esp318no93 & 1504 & Spanish & homophonic or simple + nomenclator & NS & none & 1322 \\
beatrice1505 & 1505 & Italian & not a cipher & NA & n/a & -- \\
puebla1507 & 1507 & Spanish & undetermined & DEC & complete & -- \\
maximiliana1510 & 1510 & Latin & homophonic or simple + nomenclator & DEC & partial & -- \\
vich1511 & 1511 & Spanish & syllabic nomenclator & APT & partial & 22536 \\
bonzagni1512 & 1512 & Italian & substitution, no nomenclator & DEC & partial & -- \\
vich1515 & 1515 & Spanish & homophonic or simple + nomenclator & CT & none & 3177 \\
ayala1516 & 1516 & Spanish & homophonic or simple + nomenclator & KK & complete & -- \\
schiner1516 & 1516 & Latin & undetermined & DEC & complete & -- \\
jantini1517$^\dagger$ & 1517 & Italian & undetermined & NS & none & 581 \\
r4349 & 1517 & Latin, Swedish & substitution, no nomenclator & KK & partial & 131 \\
spinelly1516 & 1517 & English & homophonic or simple + nomenclator & APT & partial & 43 \\
caprile1519 & 1519 & Italian & syllabic nomenclator & NS & none & 183 \\
caprile1520 & 1520 & Italian & substitution, no nomenclator & EPT & partial & 2247 \\
cistarelli1520 & 1520 & Italian & substitution, no nomenclator & NS & none & 708 \\
adrian1521 & 1521 & Spanish & syllabic nomenclator & EPT & complete & -- \\
lopehurtado & 1522 & Spanish & homophonic or simple + nomenclator & APT & partial & 2125 \\
poupet1522 & 1522 & French & substitution, no nomenclator & DEC & complete & -- \\
r9658 & 1522 & Spanish & homophonic or simple + nomenclator & DEC & complete & -- \\
sanchez1522 & 1522 & Spanish & homophonic or simple + nomenclator & DEC & partial & 2961 \\
frederik1523 & 1523 & German & homophonic or simple + nomenclator & DEC & complete & 308 \\
harley6345 & 1523 & English & homophonic or simple + nomenclator & NA & n/a & 78 \\
lopehurtado1523 & 1523 & Spanish & homophonic or simple + nomenclator & APT & partial & -- \\
sanchez1523 & 1523 & Spanish & homophonic or simple + nomenclator & KK & partial & -- \\
sessa1523 & 1523 & Spanish & homophonic or simple + nomenclator & EK & partial & -- \\
soria1523 & 1523 & Spanish & homophonic or simple + nomenclator & CT & partial & -- \\
moncada1524 & 1524 & Spanish & homophonic or simple + nomenclator & EPT & complete & 442 \\
sessa1524 & 1524 & Spanish & homophonic or simple + nomenclator & APT & partial & 6755 \\
ehingen1525 & 1525 & -- & undetermined & NA & n/a & -- \\
mellon29 & 1525 & -- & undetermined & NS & none & 234 \\
bayard1526 & 1526 & French & undetermined & NS & none & 56 \\
ghinucci1527 & 1526 & Latin & syllabic nomenclator & EPT & complete & 1659 \\
lee1526 & 1526 & English & homophonic or simple + nomenclator & KK & complete & 823 \\
raince$^\dagger$ & 1526 & French & homophonic or simple + nomenclator & KK & partial & 5629 \\
gilino1527$^\dagger$ & 1527 & Italian & syllabic nomenclator & EK & partial & 3111 \\
poyntz1527$^\dagger$ & 1527 & English & homophonic or simple + nomenclator & DEC & complete & 44 \\
soria1527 & 1527 & Spanish & homophonic or simple + nomenclator & DEC & complete & -- \\
vasto1527 & 1527 & Spanish & homophonic or simple + nomenclator & KK & complete & 1680 \\
dubellay1528 & 1528 & French & homophonic or simple + nomenclator & DEC & complete & -- \\
ranzo1528$^\dagger$ & 1528 & Italian & code or codebook & CT & partial & 8224 \\
visconti1528 & 1528 & Italian & homophonic or simple + nomenclator & KK & partial & 1344 \\
visconti1529 & 1528 & Italian & homophonic or simple + nomenclator & KK & partial & 1672 \\
bizozola1520 & 1529 & Italian & substitution, no nomenclator & KK & partial & 1914 \\
dubellay & 1529 & French & homophonic or simple + nomenclator & DEC & complete & 194 \\
gramont1529 & 1529 & French & homophonic or simple + nomenclator & KK & complete & 4002 \\
kaa4591b & 1529 & Latin & homophonic or simple + nomenclator & APT & partial & 6790 \\
langeac1529 & 1529 & French & substitution, no nomenclator & KK & complete & -- \\
rangone1530 & 1529 & Italian & substitution, no nomenclator & CT & partial & 264 \\
sormano$^\dagger$ & 1529 & Italian & homophonic or simple + nomenclator & KK & partial & -- \\
gramont1530 & 1530 & French & homophonic or simple + nomenclator & KK & partial & -- \\
joachim1530 & 1530 & Italian & homophonic or simple + nomenclator & CT & partial & 402 \\
isabella1531 & 1531 & Spanish & syllabic nomenclator & DEC & complete & -- \\
mai1531 & 1531 & Spanish & undetermined & NA & n/a & -- \\
muxetula1531 & 1531 & Spanish & homophonic or simple + nomenclator & DEC & complete & -- \\
augurelio1535 & 1533 & German & substitution, no nomenclator & APT & partial & 4363 \\
sperantio1534 & 1533 & German, Latin & substitution, no nomenclator & APT & partial & 7106 \\
cifuentes1534 & 1534 & Spanish & syllabic nomenclator & DEC & complete & -- \\
kaa4591 & 1534 & German, Latin & homophonic or simple + nomenclator & APT & partial & 33462 \\
marechaux1520 & 1534 & French & substitution, no nomenclator & KK & complete & 1341 \\
r9367 & 1535 & German & homophonic or simple + nomenclator & KK & partial & 534 \\
selve1535 & 1535 & French & homophonic or simple + nomenclator & KK & partial & 2524 \\
charles1536 & 1536 & Spanish & syllabic nomenclator & DEC & complete & -- \\
r4248 & 1536 & French & substitution, no nomenclator & DEC & complete & 20 \\
rome1536 & 1536 & French & homophonic or simple + nomenclator & KK & partial & 675 \\
caracciolo1537 & 1537 & Italian & syllabic nomenclator & KK & partial & -- \\
wardein1537 & 1537 & German & homophonic or simple + nomenclator & APT & partial & 4675 \\
devil1539 & 1539 & -- & undetermined & NA & n/a & 327 \\
vatican5 & 1542 & Italian & syllabic nomenclator & NS & none & 6534 \\
sadler1543 & 1543 & English & undetermined & DEC & complete & -- \\
ibrahim1548 & 1548 & Latin & homophonic or simple + nomenclator & NS & none & 572 \\
malvezzi1548 & 1548 & Latin & homophonic or simple + nomenclator & KK & complete & -- \\
guiche1551 & 1551 & French & substitution, no nomenclator & CT & partial & 431 \\
vargas1552 & 1552 & Spanish & homophonic or simple + nomenclator & DEC & complete & -- \\
grumbach1564 & 1553 & German & homophonic or simple + nomenclator & KK & partial & 692 \\
santacroce1552 & 1553 & Italian & homophonic or simple + nomenclator & DEC & complete & 3933 \\
harley1582r8499 & 1554 & English & homophonic or simple + nomenclator & EK & none & 3152 \\
renard1554 & 1554 & French & undetermined & DEC & complete & -- \\
r1212 & 1556 & Latin & homophonic or simple + nomenclator & KK & complete & 288 \\
noailles1558 & 1558 & French & undetermined & NS & none & -- \\
r1873 & 1558 & Italian & homophonic or simple + nomenclator & DEC & complete & 176 \\
burghley1559 & 1559 & English & code or codebook & DEC & complete & 12 \\
percy1559 & 1559 & -- & undetermined & NS & none & 63 \\
throck1559 & 1559 & English & homophonic or simple + nomenclator & KK & none & 729 \\
throck1559margin & 1559 & -- & code or codebook & NS & none & 222 \\
memoire1560 & 1560 & French & code or codebook & CT & partial & 1715 \\
throck1560 & 1560 & English & undetermined & DEC & complete & -- \\
throckmorton & 1560 & English & homophonic or simple + nomenclator & DEC & complete & -- \\
smith1562 & 1562 & English & homophonic or simple + nomenclator & KK & partial & -- \\
coligny1563 & 1563 & French & homophonic or simple + nomenclator & DEC & complete & -- \\
foix1563$^\dagger$ & 1563 & -- & homophonic or simple + nomenclator & EK & none & 1007 \\
lomer1563 & 1563 & English & code or codebook & DEC & complete & 7 \\
middelmore1563 & 1563 & English & code or codebook & DEC & complete & 18 \\
rennes1563 & 1563 & French & homophonic or simple + nomenclator & KK & none & 2034 \\
charlesixducroc & 1565 & French & homophonic or simple + nomenclator & DEC & complete & -- \\
toledo1565 & 1565 & Spanish & substitution, no nomenclator & APT & complete & 2087 \\
ducroc & 1567 & French & undetermined & NS & none & 166 \\
norreys1567 & 1567 & English & homophonic or simple + nomenclator & KK & partial & 7 \\
nuncio1566 & 1567 & Italian & homophonic or simple + nomenclator & KK & complete & 7817 \\
alessandrino1568 & 1568 & Italian & homophonic or simple + nomenclator & KK & partial & 10611 \\
moray & 1568 & -- & undetermined & NS & none & 134 \\
wod1568 & 1568 & -- & undetermined & NS & none & 39 \\
hamilton1569 & 1569 & Scots & homophonic or simple + nomenclator & DEC & complete & -- \\
throck1569 & 1569 & English & homophonic or simple + nomenclator & KK & complete & -- \\
norfolk1570 & 1570 & Scots & homophonic or simple + nomenclator & NS & none & -- \\
randolph1570 & 1570 & English & homophonic or simple + nomenclator & DEC & complete & -- \\
birago & 1571 & Italian & undetermined & NS & none & 483 \\
burghley1572 & 1571 & English & homophonic or simple + nomenclator & NS & none & 4 \\
r8356 & 1571 & English & homophonic or simple + nomenclator & NS & none & 11 \\
mondoucet & 1572 & French & substitution, no nomenclator & APT & partial & 7293 \\
r8358 & 1572 & English & homophonic or simple + nomenclator & NS & none & 15 \\
r8361 & 1572 & English & homophonic or simple + nomenclator & NS & none & 21 \\
walsingham1572nov & 1572 & English & undetermined & NA & n/a & -- \\
lanssac & 1573 & French & homophonic or simple + nomenclator & APT & complete & 2165 \\
ormanetto1573$^\dagger$ & 1573 & Italian & homophonic or simple + nomenclator & EK & partial & 1324 \\
r8362 & 1573 & English & homophonic or simple + nomenclator & NA & n/a & 2 \\
r8364 & 1573 & English & homophonic or simple + nomenclator & EPT & partial & 26 \\
walsingham1572 & 1573 & English & substitution, no nomenclator & NS & none & 12 \\
nevers1574 & 1574 & French, Italian & substitution, no nomenclator & DEC & partial & -- \\
philip1574 & 1574 & Spanish & homophonic or simple + nomenclator & DEC & complete & -- \\
chodkiewicz1575 & 1575 & German & substitution, no nomenclator & DEC & complete & -- \\
fulda1576 & 1576 & Latin & homophonic or simple + nomenclator & KK & partial & 270 \\
ormanetto1576$^\dagger$ & 1576 & Italian & homophonic or simple + nomenclator & EK & complete & 389 \\
ungnad1576 & 1576 & -- & homophonic or simple + nomenclator & NA & n/a & -- \\
sinzendorff1578 & 1578 & -- & undetermined & NA & n/a & -- \\
dandini1580 & 1580 & Italian & homophonic or simple + nomenclator & DEC & complete & 509 \\
dandini1580r73 & 1580 & Italian & homophonic or simple + nomenclator & DEC & partial & 1082 \\
sauli1579 & 1580 & Italian & homophonic or simple + nomenclator & DEC & complete & 2392 \\
r4157 & 1581 & French & homophonic or simple + nomenclator & KK & complete & 289 \\
walsingham1581 & 1581 & English & homophonic or simple + nomenclator & DEC & complete & 562 \\
r9319 & 1583 & De & undetermined & DEC & complete & -- \\
r9424 & 1583 & German & homophonic or simple + nomenclator & CT & partial & 913 \\
segur1583 & 1583 & French & syllabic nomenclator & NS & none & 58 \\
fleury1585 & 1584 & French & substitution, no nomenclator & NS & none & 276 \\
lodi1584 & 1584 & Spanish & homophonic or simple + nomenclator & DEC & complete & 307 \\
harley1582r8505 & 1585 & French & substitution, no nomenclator & CT & none & 3564 \\
matignon1586$^\dagger$ & 1585 & French & homophonic or simple + nomenclator & KK & partial & 17913 \\
segur & 1585 & French & syllabic nomenclator & CT & complete & 466 \\
walsingham1585 & 1585 & English & homophonic or simple + nomenclator & EPT & partial & -- \\
wotton1585 & 1585 & English & homophonic or simple + nomenclator & APT & complete & 2866 \\
bergamo1585 & 1586 & Italian & homophonic or simple + nomenclator & KK & complete & 3216 \\
clair357 & 1586 & French & homophonic or simple + nomenclator & DEC & partial & -- \\
harley1582r8500 & 1586 & English & homophonic or simple + nomenclator & APT & partial & 61 \\
harley1582r8504 & 1586 & English & homophonic or simple + nomenclator & APT & complete & 7 \\
navarre1586 & 1586 & French & syllabic nomenclator & KK & partial & 108 \\
stafford1586 & 1586 & English & homophonic or simple + nomenclator & DEC & complete & 28 \\
guise1587 & 1587 & French & homophonic or simple + nomenclator & KK & none & 2484 \\
mercoeur1587 & 1587 & French & homophonic or simple + nomenclator & KK & partial & 952 \\
morosini1588 & 1587 & Italian & homophonic or simple + nomenclator & KK & partial & 238550 \\
needham1587 & 1587 & English & substitution, no nomenclator & APT & complete & 227 \\
nevers1587 & 1587 & French & homophonic or simple + nomenclator & EK & complete & 27 \\
vieuville1587 & 1587 & French & homophonic or simple + nomenclator & EK & none & -- \\
cobham1588 & 1588 & English & homophonic or simple + nomenclator & APT & none & -- \\
harley287 & 1588 & English & homophonic or simple + nomenclator & APT & partial & -- \\
nevers1588 & 1588 & French & homophonic or simple + nomenclator & APT & complete & 476 \\
r3708 & 1588 & French & homophonic or simple + nomenclator & APT & complete & 336 \\
mendoza1589 & 1589 & Spanish & syllabic nomenclator & DEC & complete & 822 \\
morosini1589 & 1589 & Italian & syllabic nomenclator & DEC & complete & 1448 \\
orbais & 1589 & French & homophonic or simple + nomenclator & KK & partial & 91 \\
ottobon & 1589 & Italian & syllabic nomenclator & EK & none & 1528 \\
r2276 & 1589 & French & homophonic or simple + nomenclator & KK & partial & 3920 \\
bienvenut1590 & 1590 & French & homophonic or simple + nomenclator & KK & complete & 26 \\
champagne1590 & 1590 & French & homophonic or simple + nomenclator & EK & partial & 110 \\
gonzaga1590 & 1590 & Italian & homophonic or simple + nomenclator & KK & complete & 926 \\
lorraine1592$^\dagger$ & 1592 & French & homophonic or simple + nomenclator & CT & none & 1109 \\
pelissier1592 & 1592 & French & homophonic or simple + nomenclator & KK & partial & 18451 \\
ibarra1593$^\dagger$ & 1593 & Spanish & syllabic nomenclator & KK & none & 598 \\
lauriere1593$^\dagger$ & 1593 & French & syllabic nomenclator & EK & partial & 667 \\
lebel1593 & 1593 & French & homophonic or simple + nomenclator & KK & partial & 1935 \\
mauclerc1593 & 1593 & French & undetermined & DEC & complete & -- \\
mornay1593$^\dagger$ & 1593 & French & undetermined & KK & partial & 278 \\
nevers1593 & 1593 & French & homophonic or simple + nomenclator & KK & none & 395 \\
revol1593$^\dagger$ & 1593 & French & syllabic nomenclator & KK & none & 86 \\
sega1593 & 1593 & Italian & homophonic or simple + nomenclator & NS & none & -- \\
sessa1593$^\dagger$ & 1593 & Spanish & syllabic nomenclator & KK & partial & 1310 \\
stowe166 & 1593 & English & homophonic or simple + nomenclator & CT & complete & 126 \\
casimir1618 & 1594 & Dutch & undetermined & NS & none & 259 \\
joyeuse & 1594 & French & syllabic nomenclator & NS & none & 358 \\
nevers1595 & 1595 & French & homophonic or simple + nomenclator & NS & none & 753 \\
sloane3188 & 1595 & English & not a cipher & NA & n/a & -- \\
segura1596 & 1596 & Spanish & substitution, no nomenclator & CT & complete & 125 \\
bethune & 1601 & French & syllabic nomenclator & APT & partial & 2849 \\
breves1603 & 1603 & French & homophonic or simple + nomenclator & DEC & complete & -- \\
sirtori1604 & 1604 & -- & undetermined & NS & none & 747 \\
senisteros1605 & 1605 & Spanish & substitution, no nomenclator & CT & complete & 1447 \\
anhalt1610 & 1610 & German & undetermined & CT & none & 242 \\
brulart1615 & 1615 & French & homophonic or simple + nomenclator & DEC & partial & -- \\
andreae1616 & 1616 & German & substitution, no nomenclator & APT & complete & 156 \\
cocquet & 1616 & French & undetermined & NS & none & -- \\
r4817 & 1618 & German & substitution, no nomenclator & EPT & complete & 79 \\
ferdinand1619 & 1619 & German & homophonic or simple + nomenclator & CT & complete & 4181 \\
haga1620 & 1620 & Dutch & homophonic or simple + nomenclator & KK & complete & -- \\
sp106box10 & 1623 & Italian & code or codebook & KK & none & -- \\
sp106r704 & 1623 & Italian & syllabic nomenclator & NS & none & -- \\
sp106r725 & 1623 & Italian & code or codebook & NS & none & -- \\
damiata1624 & 1625 & Italian & homophonic or simple + nomenclator & DEC & complete & 361 \\
sp106r722 & 1625 & English, Italian & code or codebook & NS & none & 420 \\
warsaw & 1627 & Italian & syllabic nomenclator & CT & none & 465 \\
boswell1628 & 1628 & English & substitution, no nomenclator & KK & complete & 442 \\
kemeny1628 & 1628 & -- & undetermined & NS & none & -- \\
r1862 & 1628 & Italian & homophonic or simple + nomenclator & EPT & complete & 1280 \\
r1874 & 1628 & Italian & homophonic or simple + nomenclator & DEC & complete & 2245 \\
riksarkivet4284 & 1628 & Latin & undetermined & EK & none & 822 \\
rusdorff1628 & 1628 & German, Latin & undetermined & NS & none & -- \\
pallotto1629 & 1629 & Italian & homophonic or simple + nomenclator & DEC & complete & 118931 \\
richelieu & 1629 & French & homophonic or simple + nomenclator & APT & complete & 421 \\
sg1629 & 1629 & -- & undetermined & NA & n/a & -- \\
riksarkivet1628 & 1631 & German & homophonic or simple + nomenclator & APT & partial & 1130 \\
ceva1632 & 1632 & Italian & syllabic nomenclator & DEC & complete & 2400 \\
oxenstierna1633 & 1633 & Swedish & homophonic or simple + nomenclator & NS & none & -- \\
ferdinand1634 & 1634 & German & homophonic or simple + nomenclator & NS & none & -- \\
ferdinand3 & 1634 & Latin & syllabic nomenclator & APT & partial & 1719 \\
ormonde & 1634 & English & homophonic or simple + nomenclator & APT & partial & 59 \\
harley286 & 1635 & Latin & substitution, no nomenclator & CT & partial & 281 \\
koerler1675 & 1635 & German & undetermined & NA & n/a & -- \\
marburg1635 & 1635 & German & substitution, no nomenclator & CT & partial & 590 \\
emery1636 & 1636 & French & undetermined & NS & none & -- \\
rohan1636 & 1636 & French & undetermined & NS & none & -- \\
heusner1637 & 1637 & German & homophonic or simple + nomenclator & KK & complete & 259 \\
malsburg1637$^\dagger$ & 1637 & French, German & transposition, polygraphic, polyalphabetic & CT & partial & 13487 \\
kurtz1639 & 1639 & German & substitution, no nomenclator & APT & partial & 22913 \\
baner1640 & 1640 & German & homophonic or simple + nomenclator & APT & partial & 288 \\
harley7001 & 1641 & English & homophonic or simple + nomenclator & EK & partial & 254 \\
joao1642 & 1642 & Portuguese & homophonic or simple + nomenclator & APT & partial & 1817 \\
boswell & 1643 & English & homophonic or simple + nomenclator & KK & complete & 457 \\
forster & 1644 & French & substitution, no nomenclator & KK & complete & 207 \\
lucca & 1644 & Italian & substitution, no nomenclator & EPT & partial & 231 \\
goring1645 & 1645 & English & substitution, no nomenclator & KK & complete & 338 \\
hm1645 & 1645 & English & homophonic or simple + nomenclator & DEC & complete & 3373 \\
rupert & 1645 & English & homophonic or simple + nomenclator & KK & partial & 1412 \\
rupert1645 & 1645 & English & homophonic or simple + nomenclator & EK & complete & 304 \\
sp106r787 & 1645 & French & homophonic or simple + nomenclator & DEC & complete & 1569 \\
sp106r931 & 1645 & English & code or codebook & NS & none & -- \\
transylvania1645 & 1645 & Hungarian & undetermined & CT & partial & 4785 \\
sp106r924 & 1646 & English & homophonic or simple + nomenclator & DEC & complete & 2012 \\
sp106r929 & 1646 & English & homophonic or simple + nomenclator & KK & complete & 348 \\
charlesi & 1648 & English & undetermined & NS & none & 200 \\
mercy1648$^\dagger$ & 1648 & Spanish & homophonic or simple + nomenclator & CT & partial & 529 \\
r4941 & 1648 & -- & undetermined & NS & none & -- \\
conti1649 & 1649 & French & substitution, no nomenclator & KK & complete & 1692 \\
laigue1649 & 1649 & French & syllabic nomenclator & APT & complete & 3367 \\
hamilton & 1650 & English & homophonic or simple + nomenclator & NA & n/a & 102 \\
wallis1653 & 1651 & English, French & syllabic nomenclator & DEC & complete & 480 \\
conley1652 & 1652 & English & homophonic or simple + nomenclator & NS & none & 23 \\
urquhart & 1652 & English & code or codebook & KK & partial & 349 \\
bordeaux & 1653 & French & syllabic nomenclator & EK & complete & 813 \\
thurloe & 1653 & -- & homophonic or simple + nomenclator & NS & none & 136 \\
comenius1655$^\dagger$ & 1655 & Latin & substitution, no nomenclator & DEC & complete & 88 \\
copinger1656 & 1656 & English & undetermined & NS & none & 16 \\
dugard1656 & 1656 & -- & undetermined & NS & none & 27 \\
waddall1656 & 1656 & English & homophonic or simple + nomenclator & NS & none & 50 \\
letellier & 1657 & French & undetermined & NS & none & 202 \\
rome1658 & 1658 & Italian & undetermined & NA & n/a & -- \\
hyde & 1659 & -- & not a cipher & NA & n/a & 41 \\
miranda1660 & 1660 & French, Italian, Latin, Portuguese, Spanish & undetermined & NA & n/a & -- \\
r688 & 1664 & Serbian & undetermined & NA & n/a & 240 \\
wesselenyi1664 & 1664 & Latin & homophonic or simple + nomenclator & CT & complete & 244 \\
erbach1653 & 1666 & German, Latin & substitution, no nomenclator & APT & complete & 610 \\
perwich & 1670 & English & transposition, polygraphic, polyalphabetic & DEC & complete & 494 \\
hesse1672 & 1672 & German & syllabic nomenclator & EK & partial & 40 \\
hernannunez1674 & 1674 & Spanish & syllabic nomenclator & KK & complete & 1716 \\
biermann1675$^\dagger$ & 1675 & French & syllabic nomenclator & EK & partial & 263 \\
colbert & 1675 & French & undetermined & NS & none & 106 \\
balbases1677 & 1677 & Spanish & syllabic nomenclator & DEC & partial & -- \\
carpio1677 & 1677 & Spanish & syllabic nomenclator & APT & complete & 3504 \\
ronquillo1676 & 1677 & Spanish & syllabic nomenclator & DEC & partial & 23463 \\
ormond1678 & 1678 & English & code or codebook & NS & none & 20 \\
r1871 & 1680 & Italian & homophonic or simple + nomenclator & DEC & complete & 610 \\
szembek1689 & 1689 & Latin & homophonic or simple + nomenclator & DEC & complete & 439 \\
chaulnes & 1690 & French & syllabic nomenclator & NS & none & 311 \\
feuquieres & 1690 & French & code or codebook & EPT & partial & 601 \\
catinat1691 & 1691 & French & code or codebook & KK & complete & 12439 \\
ilten1697 & 1697 & French & code or codebook & DEC & complete & 74 \\
yard1699 & 1699 & English & code or codebook & KK & complete & 368 \\
catinat1702 & 1702 & French & code or codebook & NS & none & 591 \\
prussia1702 & 1702 & French & undetermined & NA & n/a & 56 \\
stepney & 1702 & English & code or codebook & NS & none & 24 \\
ilten1705 & 1705 & French, German & code or codebook & DEC & complete & 418 \\
bay1706 & 1706 & Hungarian & syllabic nomenclator & EK & partial & 414 \\
papai1706 & 1706 & Hungarian & syllabic nomenclator & KK & partial & 16861 \\
rakoczi1704 & 1706 & French & syllabic nomenclator & EK & partial & 2875 \\
schenck1706 & 1706 & German & homophonic or simple + nomenclator & CT & partial & 1181 \\
vetesi1705 & 1706 & Hungarian & syllabic nomenclator & EK & complete & -- \\
henter1707 & 1707 & Hungarian & homophonic or simple + nomenclator & DEC & complete & -- \\
r7475 & 1707 & Polish & homophonic or simple + nomenclator & DEC & complete & 2444 \\
rakoczi1707 & 1707 & French & syllabic nomenclator & EK & complete & 1108 \\
sieniawski1708 & 1709 & Polish & homophonic or simple + nomenclator & DEC & complete & -- \\
turculet1709 & 1709 & Polish & homophonic or simple + nomenclator & DEC & complete & 877 \\
desmarets1710 & 1710 & French & homophonic or simple + nomenclator & KK & partial & 471 \\
geertruidenberg & 1710 & French & code or codebook & NS & none & 690 \\
papai1710 & 1710 & Hungarian & undetermined & NS & none & 371 \\
villars1710 & 1710 & French & code or codebook & NS & none & 173 \\
rakoczi1711 & 1711 & Latin & syllabic nomenclator & DEC & partial & 2531 \\
anne1711 & 1712 & English & code or codebook & DEC & complete & 414 \\
lagau1712 & 1712 & French & undetermined & NA & n/a & -- \\
chomentowski1713 & 1713 & Polish & syllabic nomenclator & DEC & partial & 2246 \\
mniszech1714 & 1714 & Latin, Polish & undetermined & NS & none & 232 \\
stenbock1713 & 1714 & -- & undetermined & NA & n/a & -- \\
potocka1714 & 1715 & French, Polish & homophonic or simple + nomenclator & CT & partial & 886 \\
goertz1717 & 1717 & German & homophonic or simple + nomenclator & NS & none & 809 \\
duran1718 & 1718 & -- & code or codebook & NS & none & 317 \\
nunzio1718 & 1718 & Italian & homophonic or simple + nomenclator & KK & complete & 34243 \\
bl32305 & 1719 & French & code or codebook & NS & none & 354 \\
windischgraetz1720 & 1720 & German & syllabic nomenclator & KK & complete & 223 \\
windischgraetz1720b & 1720 & German & homophonic or simple + nomenclator & KK & complete & 192 \\
colonia1721 & 1721 & Italian & homophonic or simple + nomenclator & KK & complete & 189 \\
windischgraetz1721 & 1721 & German & homophonic or simple + nomenclator & KK & partial & 181 \\
palm1727 & 1727 & German & homophonic or simple + nomenclator & NS & none & 185 \\
visconti1727 & 1727 & Italian & homophonic or simple + nomenclator & NA & n/a & 193 \\
r7707 & 1728 & -- & undetermined & NS & none & 205 \\
portstemarie1729 & 1729 & French & code or codebook & APT & complete & 141 \\
sp106r597 & 1731 & French & homophonic or simple + nomenclator & NA & n/a & 43 \\
august1733 & 1733 & -- & undetermined & NA & n/a & -- \\
mecklenburg1734 & 1734 & German & code or codebook & NS & none & 598 \\
nunzio1736 & 1736 & Italian & homophonic or simple + nomenclator & KK & partial & 7386 \\
carre1742 & 1742 & -- & code or codebook & NS & none & 1323 \\
ilten1743 & 1743 & French & code or codebook & EK & complete & 223 \\
esterhazy1744 & 1744 & French & code or codebook & NS & none & 176 \\
waldeck1744 & 1744 & German & homophonic or simple + nomenclator & NS & none & 717 \\
sp95box136 & 1748 & -- & undetermined & NA & n/a & -- \\
michell1751$^\dagger$ & 1751 & French & code or codebook & DEC & complete & 550 \\
hellen1752 & 1752 & French & code or codebook & NS & none & 2651 \\
jones1753 & 1753 & English & not a cipher & NA & n/a & -- \\
kauderbach1754 & 1754 & French & homophonic or simple + nomenclator & CT & partial & 8512 \\
affry1757 & 1757 & French & homophonic or simple + nomenclator & APT & partial & 8980 \\
affry1757c & 1757 & French & homophonic or simple + nomenclator & NS & none & 761 \\
acciaiuoli1757 & 1758 & Italian & homophonic or simple + nomenclator & KK & complete & 8326 \\
eichel1758 & 1758 & -- & homophonic or simple + nomenclator & NA & n/a & 139 \\
swieten1757 & 1758 & French & homophonic or simple + nomenclator & DEC & complete & 89 \\
orvilliers1759 & 1759 & -- & undetermined & NA & n/a & -- \\
swieten1759 & 1759 & French & homophonic or simple + nomenclator & CT & complete & 553 \\
mercier1761 & 1761 & French & code or codebook & APT & partial & 818 \\
r1875 & 1761 & Italian & substitution, no nomenclator & DEC & complete & 1565 \\
duben1764$^\dagger$ & 1764 & French & undetermined & NS & none & 4340 \\
lucini1767 & 1767 & Italian & homophonic or simple + nomenclator & KK & complete & 152 \\
smallwood1777 & 1777 & English & not a cipher & NA & n/a & 5 \\
martinique1778 & 1778 & French & code or codebook & APT & partial & 285 \\
destaing & 1779 & French & code or codebook & NS & none & 217 \\
pcc1780 & 1780 & English & code or codebook & DEC & complete & -- \\
balfour1781 & 1781 & English & transposition, polygraphic, polyalphabetic & DEC & complete & 6337 \\
livingston1781 & 1781 & English & homophonic or simple + nomenclator & DEC & complete & 17 \\
deswart1782 & 1782 & Dutch & code or codebook & EK & partial & 9603 \\
schoning1783 & 1783 & French & code or codebook & APT & complete & 4100 \\
rechteren1785 & 1785 & Dutch & homophonic or simple + nomenclator & DEC & complete & 424 \\
jefferson1787 & 1786 & English & code or codebook & APT & partial & 218 \\
vanreede1787 & 1787 & French & code or codebook & KK & complete & 471 \\
deboli1790 & 1788 & Polish & code or codebook & DEC & complete & 211 \\
dedem1788$^\dagger$ & 1788 & Dutch & code or codebook & DEC & partial & 941 \\
legation1788 & 1788 & English & code or codebook & DEC & complete & 84 \\
deboli1789 & 1789 & Polish & syllabic nomenclator & APT & partial & 29465 \\
reede1789 & 1789 & French & code or codebook & NS & none & 4296 \\
deboli1791 & 1791 & Polish & code or codebook & DEC & partial & 3082 \\
jablonowski1791 & 1791 & French & undetermined & NA & n/a & -- \\
vanreede1792 & 1792 & French & code or codebook & KK & complete & 83 \\
vanreede1793$^\dagger$ & 1793 & Dutch & code or codebook & NS & none & 696 \\
r1892 & 1795 & Dutch, French & substitution, no nomenclator & CT & partial & 2067 \\
r2236 & 1795 & Dutch & undetermined & NS & none & -- \\
r2242 & 1795 & Dutch & substitution, no nomenclator & EK & partial & 1261 \\
wilhelmina1795 & 1795 & Dutch, French & code or codebook & NS & none & -- \\
wilkinson1789 & 1795 & English & code or codebook & NS & none & 3408 \\
hereditary1796 & 1796 & French & code or codebook & APT & complete & 881 \\
r2234 & 1796 & French & transposition, polygraphic, polyalphabetic & KK & complete & 1321 \\
reede1796 & 1796 & French & code or codebook & NS & none & 1512 \\
yvoy1797 & 1797 & Dutch & code or codebook & DEC & complete & 3487 \\
astrea1799$^\dagger$ & 1798 & French & undetermined & NA & n/a & -- \\
jones1802 & 1798 & English & code or codebook & EPT & partial & 1049 \\
mellon136 & 1798 & German & substitution, no nomenclator & KK & complete & 1670 \\
dedem1799 & 1799 & -- & code or codebook & NS & none & -- \\
king1801 & 1801 & English & undetermined & DEC & complete & -- \\
mountflorence1801 & 1801 & English & code or codebook & CT & complete & 39 \\
murray1801 & 1801 & English & code or codebook & APT & complete & 6 \\
r1944 & 1801 & Dutch & code or codebook & DEC & complete & 110 \\
r2232 & 1801 & French & substitution, no nomenclator & KK & complete & 1065 \\
r1942 & 1803 & Dutch & code or codebook & NS & none & 286 \\
toulon1803 & 1803 & Dutch & code or codebook & NS & none & 110 \\
fagel1804 & 1804 & French & homophonic or simple + nomenclator & NS & none & 128 \\
graham1806 & 1806 & English & code or codebook & NS & none & -- \\
erving1807 & 1807 & English & code or codebook & DEC & complete & 581 \\
armstrong & 1808 & English & code or codebook & EPT & complete & 49 \\
spaen1808 & 1808 & -- & code or codebook & NS & none & 303 \\
roell1809 & 1809 & -- & code or codebook & NS & none & 2585 \\
jqa & 1812 & English & code or codebook & APT & partial & 390 \\
poinsett1813 & 1813 & English & code or codebook & NS & none & 560 \\
castelcicala1816 & 1816 & Italian & code or codebook & DEC & partial & 9428 \\
loeschnen1816 & 1816 & German & not a cipher & NA & n/a & -- \\
hesse1824 & 1824 & German & transposition, polygraphic, polyalphabetic & APT & complete & 164 \\
zeschau1841 & 1841 & French, German & syllabic nomenclator & NS & none & 1896 \\
soglia1848 & 1848 & Italian & code or codebook & CT & complete & 955 \\
civilwar1862 & 1862 & English & undetermined & NA & n/a & -- \\
eckert1862 & 1862 & English & transposition, polygraphic, polyalphabetic & KK & none & 100484 \\
milroy & 1862 & English & transposition, polygraphic, polyalphabetic & KK & complete & 79 \\
union1864 & 1864 & English & transposition, polygraphic, polyalphabetic & DEC & n/a & 59 \\
army1865 & 1865 & English & undetermined & NA & n/a & -- \\
vienna1873 & 1870 & English & code or codebook & NS & none & -- \\
augur1874 & 1874 & English & transposition, polygraphic, polyalphabetic & CT & partial & 246 \\
catokwacopa & 1875 & English & transposition, polygraphic, polyalphabetic & NS & none & 441 \\
bell1900 & 1900 & English & code or codebook & KK & complete & 7 \\
gupta1917 & 1916 & English & undetermined & NA & n/a & -- \\
huangxing$^\dagger$ & 1916 & Chinese & transposition, polygraphic, polyalphabetic & KK & complete & 275 \\
sunyatsen & 1916 & Chinese & transposition, polygraphic, polyalphabetic & CT & complete & 177 \\
censorship1917 & 1917 & -- & undetermined & NA & n/a & -- \\
lima1917 & 1917 & German & undetermined & NA & n/a & -- \\
tampico1917 & 1917 & -- & undetermined & NA & n/a & -- \\
adfgvx & 1918 & German & transposition, polygraphic, polyalphabetic & KK & partial & 3966 \\
ludendorff1918 & 1918 & -- & undetermined & NS & none & 328 \\
stockholm1918 & 1918 & -- & code or codebook & NS & none & 938 \\
karakhan1921 & 1921 & -- & undetermined & NS & none & 181 \\
kopp1921 & 1921 & -- & undetermined & NS & none & 717 \\
lejava1921 & 1921 & Ru & transposition, polygraphic, polyalphabetic & CT & complete & 167 \\
stomoniakov1921 & 1921 & -- & undetermined & NS & none & 1042 \\
voikov1921 & 1921 & Ru & transposition, polygraphic, polyalphabetic & CT & complete & 615 \\
voikov1921b & 1921 & Ru & transposition, polygraphic, polyalphabetic & NS & none & 398 \\
roosevelt & 1935 & -- & not a cipher & NA & n/a & 52 \\
blume & 1937 & Spanish & transposition, polygraphic, polyalphabetic & NS & none & 775 \\
zhongshan & 1938 & Chinese & code or codebook & NA & n/a & 7 \\
dagapeyeff & 1939 & -- & undetermined & NS & none & 196 \\
orpo1942 & 1942 & German & transposition, polygraphic, polyalphabetic & NS & none & 968 \\
orpo1942jun & 1942 & German & transposition, polygraphic, polyalphabetic & DEC & complete & 223 \\
powgerman1942 & 1942 & German & undetermined & NS & none & -- \\
censorship & 1943 & German & substitution, no nomenclator & NS & none & -- \\
abwehr & 1944 & -- & undetermined & NS & none & 924 \\
ccr1944 & 1944 & -- & undetermined & NA & n/a & -- \\
buenosaires1945 & 1945 & German & undetermined & NA & n/a & -- \\
enigma & 1945 & German & transposition, polygraphic, polyalphabetic & NS & none & 97 \\
blancmesnil & -- & French & undetermined & NA & n/a & -- \\
burgo1527 & -- & Spanish & not a cipher & NA & n/a & -- \\
debosnys & -- & -- & undetermined & NS & none & 975 \\
egmond & -- & French & substitution, no nomenclator & CT & complete & 774 \\
fr3029 & -- & French & undetermined & KK & complete & -- \\
fr3029f134 & -- & French & homophonic or simple + nomenclator & KK & complete & 926 \\
fr3092 & -- & French & homophonic or simple + nomenclator & KK & complete & 4521 \\
mercoeur1586 & -- & French & homophonic or simple + nomenclator & NS & none & 217 \\
mid4131829 & -- & English & undetermined & NA & n/a & -- \\
r1818 & -- & -- & undetermined & NS & none & 355 \\
r8870 & -- & -- & code or codebook & NS & none & 901 \\
r8871 & -- & -- & undetermined & NS & none & 77 \\
r9416 & -- & German & substitution, no nomenclator & CT & partial & 3313 \\
r9586 & -- & Italian & code or codebook & NS & none & 969 \\
siena10 & -- & Italian & substitution, no nomenclator & DEC & complete & 25 \\
siena19 & -- & Italian & undetermined & NS & none & 118 \\
siena21 & -- & Italian & undetermined & NS & none & 107 \\
siena7 & -- & Italian & undetermined & NS & none & 363 \\
sp106r927 & -- & French & substitution, no nomenclator & CT & partial & 1348 \\
sp53 & -- & -- & undetermined & NS & none & 1235 \\
\end{longtable}}

{\scriptsize
\begin{longtable}{@{}llrlllr@{}}
\caption{\label{tab:allcontrib} The contributed targets, by date of the earliest document, with the credited solver (Table~\ref{tab:contributors}). Methods as in Table~\ref{tab:all}.}\\
\toprule
target & solver & year & language & method & extent & tokens \\
\midrule
\endfirsthead
\toprule
target & solver & year & language & method & extent & tokens \\
\midrule
\endhead
\bottomrule
\endlastfoot
siena12 & \Who{C1}{George Lasry} & 1413 & Italian, Latin & CT & partial & 865 \\
siena2023 & \Who{C1}{George Lasry} & 1413 & Italian, Latin & CT & complete & 4254 \\
siena1431 & \Who{C1}{George Lasry} & 1431 & Italian & CT & partial & 3694 \\
it1583 & \Who{C1}{George Lasry} & 1446 & Italian & CT & partial & 827 \\
siena9 & \Who{C1}{George Lasry} & 1448 & Italian & APT & partial & 53 \\
tranchedini1452 & \Who{C1}{George Lasry} & 1452 & Italian & DEC & partial & 852 \\
ferdinand1473 & \Who{C4}{Michael Brackx} & 1473 & Spanish & APT & partial & 511 \\
siena11 & \Who{C1}{George Lasry} & 1478 & Italian & CT & partial & 393 \\
buda1490 & \Who{C12}{Harrison Bornstein} & 1490 & Latin & EPT & partial & 3199 \\
puebla1495 & \Who{C3}{WAINUT} & 1495 & Spanish & KK & partial & -- \\
puebla1496 & \Who{C3}{WAINUT} & 1496 & Spanish & DEC & complete & -- \\
esp318no94 & \Who{C4}{Michael Brackx} & 1503 & Spanish & CT & none & 3554 \\
puebla1506 & \Who{C3}{WAINUT} & 1506 & Spanish & KK & partial & 373 \\
dupuy468 & \Who{C1}{George Lasry} & 1519 & Latin & APT & partial & 684 \\
adrian1520 & \Who{C2}{Brice Dudout} & 1520 & Spanish & EPT & complete & 113 \\
lopehurtado1520 & \Who{C2}{Brice Dudout} & 1520 & Spanish & CT & complete & 3736 \\
nino1522 & \Who{C1}{George Lasry} & 1522 & Spanish & CT & complete & 260 \\
salviati1525 & \Who{C1}{George Lasry} & 1525 & Italian & CT & partial & 838 \\
bilia1527 & \Who{C6}{Harumi-Matsuba} & 1527 & Italian & KK & partial & 623 \\
siena17 & \Who{C1}{George Lasry} & 1528 & Italian & CT & partial & 282 \\
turenne1528 & \Who{C1}{George Lasry} & 1528 & French & APT & partial & 785 \\
ferducci1532 & \Who{C2}{Brice Dudout} & 1532 & Italian & APT & complete & 8 \\
huyson1547 & \Who{C2}{Brice Dudout} & 1547 & Latin & EK & complete & 8517 \\
siena15 & \Who{C1}{George Lasry} & 1547 & Italian, Latin & CT & none & 240 \\
siguenza1554 & \Who{C4}{Michael Brackx} & 1554 & Spanish & APT & partial & 887 \\
dubellay1555 & \Who{C1}{George Lasry} & 1555 & French & DEC & partial & 193 \\
renard1556 & \Who{C1}{George Lasry} & 1556 & French & DEC & partial & 456 \\
seure1558 & \Who{C9}{Angus Robinson} & 1558 & French & APT & partial & 9405 \\
paleologue1562 & \Who{C1}{George Lasry} & 1562 & Italian & CT & partial & 5 \\
biglia1565 & \Who{C2}{Brice Dudout} & 1565 & Italian & CT & partial & 381 \\
fiesque1567 & \Who{C6}{Harumi-Matsuba} & 1567 & French & KK & partial & 29707 \\
fogliani1568 & \Who{C1}{George Lasry} & 1568 & Italian & CT & complete & 103 \\
birago1571 & \Who{C1}{George Lasry} & 1571 & Italian & EK & partial & 923 \\
vulcob1573 & \Who{C1}{George Lasry} & 1573 & French & APT & complete & 754 \\
ancel1577 & \Who{C6}{Harumi-Matsuba} & 1577 & French & KK & partial & 4385 \\
labbe1582 & \Who{C9}{Angus Robinson} & 1577 & French & CT & complete & 576 \\
perez1578 & \Who{C2}{Brice Dudout} & 1578 & Spanish & KK & complete & 153 \\
entragues1581 & \Who{C1}{George Lasry} & 1581 & French & KK & partial & 87 \\
scipion1585 & \Who{C1}{George Lasry} & 1585 & Italian & APT & partial & 3078 \\
camus1586 & \Who{C4}{Michael Brackx} & 1586 & French & DEC & partial & 2452 \\
malicorne1586 & \Who{C4}{Michael Brackx} & 1586 & French & KK & partial & 1144 \\
matignon1586c3 & \Who{C8}{alejandrozarco}, \Who{C10}{Brett Burford} & 1586 & French & KK & partial & 4079 \\
matignon1586tuple & \Who{C4}{Michael Brackx} & 1586 & French & DEC & partial & 2152 \\
mayenne1586c2 & \Who{C10}{Brett Burford}, \Who{C6}{Harumi-Matsuba} & 1586 & French & KK & partial & 5293 \\
saintluc1586 & \Who{C1}{George Lasry} & 1586 & French & CT & complete & 1315 \\
saintluc1586b & \Who{C4}{Michael Brackx} & 1586 & French & KK & none & 2428 \\
lucinge1588 & \Who{C1}{George Lasry} & 1588 & French & CT & partial & 14608 \\
dampierre1589 & \Who{C1}{George Lasry} & 1589 & French & APT & complete & 67 \\
fr4696num & \Who{C1}{George Lasry} & 1589 & Italian & APT & partial & 1744 \\
meryvic1589 & \Who{C1}{George Lasry} & 1589 & French & CT & partial & 545 \\
nevers1589 & \Who{C5}{Arya Sanketbhai Patel} & 1589 & French & NA & n/a & -- \\
longlee1590 & \Who{C2}{Brice Dudout} & 1590 & French & EPT & partial & 2044 \\
rondinelli1590 & \Who{C1}{George Lasry} & 1590 & Italian & EK & complete & 656 \\
sega1590 & \Who{C2}{Brice Dudout} & 1590 & Italian & KK & complete & 530 \\
sega1591 & \Who{C1}{George Lasry} & 1591 & Italian & KK & complete & 100 \\
dinteville1592 & \Who{C8}{alejandrozarco} & 1592 & French & APT & partial & 702 \\
maisse1592 & \Who{C5}{Arya Sanketbhai Patel} & 1592 & French & DEC & complete & 2293 \\
desportes1593 & \Who{C13}{Paolo Rosson} & 1593 & French & KK & complete & 6699 \\
maisse1593 & \Who{C5}{Arya Sanketbhai Patel}, \Who{C3}{WAINUT} & 1593 & French & DEC & partial & -- \\
gondi1594 & \Who{C1}{George Lasry} & 1594 & Italian & APT & partial & 1469 \\
nevers1594gondi & \Who{C1}{George Lasry} & 1594 & French & EK & partial & 245 \\
blasco1597 & \Who{C1}{George Lasry} & 1597 & Spanish & CT & partial & 107 \\
dupuy155 & \Who{C1}{George Lasry} & 1597 & Italian & APT & partial & 1952 \\
sessa1597 & \Who{C1}{George Lasry} & 1597 & Spanish & KK & partial & 379 \\
aerssen1601 & \Who{C7}{Feyseel Nur} & 1601 & French & APT & partial & 72 \\
hesse1603 & \Who{C5}{Arya Sanketbhai Patel} & 1602 & French & KK & complete & 4100 \\
berny1607 & \Who{C1}{George Lasry} & 1607 & French & APT & complete & 266 \\
jeannin1607 & \Who{C1}{George Lasry} & 1607 & French & APT & partial & 1145 \\
laboderie1608 & \Who{C1}{George Lasry} & 1608 & French & APT & partial & 1023 \\
bianchi & \Who{C1}{George Lasry} & 1610 & Italian & KK & complete & 256 \\
breves1610 & \Who{C2}{Brice Dudout}, \Who{C1}{George Lasry} & 1610 & French & EK & partial & 864 \\
herbault1625 & \Who{C1}{George Lasry} & 1625 & French & APT & partial & 375 \\
rusdorff1625 & \Who{C1}{George Lasry} & 1625 & Latin & CT & complete & 438 \\
herbault1626 & \Who{C5}{Arya Sanketbhai Patel} & 1626 & French & DEC & complete & -- \\
bethune1629 & \Who{C1}{George Lasry} & 1629 & French & APT & partial & 177 \\
santacruz1630 & \Who{C2}{Brice Dudout} & 1630 & Spanish & EPT & complete & 2733 \\
sabran1633 & \Who{C1}{George Lasry} & 1631 & French & CT & partial & 3506 \\
chatillon1635 & \Who{C1}{George Lasry} & 1635 & French & APT & complete & 666 \\
louis1635 & \Who{C1}{George Lasry} & 1635 & French & APT & partial & 967 \\
vizani1637 & \Who{C15}{Satoru} & 1637 & Italian & EK & partial & 190 \\
genes1638 & \Who{C2}{Brice Dudout} & 1638 & Spanish & EPT & complete & 652 \\
genes1643 & \Who{C2}{Brice Dudout} & 1643 & Spanish & EPT & complete & 634 \\
albornoz1648 & \Who{C2}{Brice Dudout} & 1644 & Spanish & EPT & complete & 1696 \\
letellier1644 & \Who{C6}{Harumi-Matsuba} & 1644 & French & KK & partial & 626 \\
brienne1067 & \Who{C1}{George Lasry} & 1646 & French & APT & complete & 350 \\
bouillon1650 & \Who{C1}{George Lasry} & 1650 & French & APT & complete & 949 \\
clair460alpha & \Who{C1}{George Lasry} & 1650 & French & APT & complete & 888 \\
bagno1652 & \Who{C1}{George Lasry} & 1652 & Italian & CT & complete & 1796 \\
brienne5160 & \Who{C1}{George Lasry} & 1653 & French & APT & complete & 1356 \\
sande1660 & \Who{C3}{WAINUT} & 1657 & Portuguese & CT & partial & 1336 \\
asang1662 & \Who{C2}{Brice Dudout} & 1662 & Polish & CT & complete & 320 \\
wasowicz1662 & \Who{C2}{Brice Dudout} & 1662 & Polish & KK & complete & 1587 \\
vilaflor1663 & \Who{C2}{Brice Dudout} & 1663 & Portuguese & CT & complete & 1953 \\
portugal1667 & \Who{C1}{George Lasry} & 1667 & Portuguese & DEC & complete & 257 \\
beuningen1676 & \Who{C7}{Feyseel Nur} & 1676 & Dutch & APT & partial & 278 \\
feuquieres1678 & \Who{C1}{George Lasry} & 1678 & French & APT & partial & 1920 \\
duquesne1681 & \Who{C1}{George Lasry} & 1681 & French & APT & complete & 751 \\
battier1686 & \Who{C7}{Feyseel Nur}, \Who{C3}{WAINUT} & 1686 & Dutch & KK & partial & 2560 \\
vergier1696 & \Who{C1}{George Lasry} & 1696 & French & APT & complete & 1212 \\
bossuet & \Who{C1}{George Lasry} & 1697 & French & KK & partial & 642 \\
stade1710 & \Who{C2}{Brice Dudout} & 1710 & Swedish & APT & complete & 640 \\
clair1225 & \Who{C1}{George Lasry} & 1714 & French & APT & complete & 495 \\
partyet1721 & \Who{C2}{Brice Dudout} & 1721 & French & KK & complete & 77 \\
villardeau & \Who{C1}{George Lasry} & 1725 & French & APT & complete & 1248 \\
blathwayt1729 & \Who{C2}{Brice Dudout} & 1729 & French & KK & complete & 141 \\
newcastle1729 & \Who{C14}{johonn} & 1729 & French & APT & none & 141 \\
villeneuve1737 & \Who{C1}{George Lasry} & 1737 & French & CT & complete & 815 \\
champigny1744 & \Who{C2}{Brice Dudout} & 1744 & French & APT & complete & 381 \\
r1955 & \Who{C3}{WAINUT} & 1751 & French & EK & none & 362 \\
mondesir1769 & \Who{C1}{George Lasry} & 1769 & French & CT & partial & 4923 \\
bouille1778 & \Who{C2}{Brice Dudout} & 1778 & French & KK & complete & 759 \\
lovell1781 & \Who{C3}{WAINUT} & 1781 & English & KK & complete & 57 \\
dana1783 & \Who{C3}{WAINUT} & 1783 & English & CT & partial & 212 \\
damas1787 & \Who{C2}{Brice Dudout} & 1787 & French & KK & complete & 463 \\
deboli1788 & \Who{C2}{Brice Dudout} & 1788 & Polish & APT & partial & 943 \\
erving1807aug & \Who{C7}{Feyseel Nur} & 1807 & English & EPT & complete & 1145 \\
marmont1812 & \Who{C2}{Brice Dudout} & 1812 & French & KK & partial & 443 \\
napoleon & \Who{C2}{Brice Dudout} & 1812 & French & KK & complete & 325 \\
crawford1813 & \Who{C3}{WAINUT} & 1813 & English & KK & partial & 1368 \\
morris1813 & \Who{C3}{WAINUT} & 1814 & English & EK & partial & 554 \\
revenga1822 & \Who{C3}{WAINUT} & 1822 & Spanish & CT & partial & 2294 \\
appleton1823 & \Who{C3}{WAINUT} & 1823 & English & KK & partial & 18 \\
everett1825 & \Who{C3}{WAINUT}, \Who{C2}{Brice Dudout} & 1825 & English & KK & partial & 355 \\
ellis1839 & \Who{C3}{WAINUT}, \Who{C4}{Michael Brackx} & 1836 & English & DEC & complete & 609 \\
eugenie1859 & \Who{C1}{George Lasry} & 1859 & French & CT & partial & 440 \\
greenhow1861 & \Who{C3}{WAINUT} & 1861 & English & EK & partial & 200 \\
jordan1861 & \Who{C3}{WAINUT} & 1861 & English & DEC & complete & 89 \\
johnston1863 & \Who{C3}{WAINUT} & 1863 & English & DEC & complete & 31 \\
taylor1865 & \Who{C3}{WAINUT} & 1865 & English & KK & partial & 1680 \\
army1902 & \Who{C6}{Harumi-Matsuba} & 1902 & Japanese & APT & partial & 73 \\
carlos1903 & \Who{C3}{WAINUT} & 1903 & French, Portuguese & DEC & partial & 74 \\
hintze1903 & \Who{C3}{WAINUT} & 1903 & Portuguese & DEC & partial & 54 \\
sabugosa1910 & \Who{C3}{WAINUT} & 1903 & Portuguese & DEC & partial & 106 \\
lima1909 & \Who{C3}{WAINUT} & 1909 & Portuguese & DEC & complete & 409 \\
manuel1910 & \Who{C3}{WAINUT} & 1910 & Portuguese & NS & none & 93 \\
ctcplain1913 & \Who{C2}{Brice Dudout} & 1913 & Chinese & KK & complete & 765 \\
rubio1913 & \Who{C2}{Brice Dudout} & 1913 & Spanish & KK & complete & 223 \\
hongkong1914 & \Who{C2}{Brice Dudout} & 1914 & Chinese & CT & complete & 44 \\
hsuweilin1914 & \Who{C2}{Brice Dudout} & 1914 & Chinese & CT & complete & 19 \\
laisaosang1915 & \Who{C2}{Brice Dudout} & 1914 & Chinese & CT & complete & 77 \\
newyork1915 & \Who{C2}{Brice Dudout} & 1914 & Chinese & CT & complete & 109 \\
rosamonde1915 & \Who{C2}{Brice Dudout} & 1915 & English & CT & partial & -- \\
thirtieth1915 & \Who{C2}{Brice Dudout} & 1915 & Chinese & CT & complete & 263 \\
numeric1916 & \Who{C2}{Brice Dudout} & 1916 & Chinese & CT & partial & 2100 \\
seventeenth1916 & \Who{C2}{Brice Dudout} & 1916 & Chinese & CT & complete & 97 \\
sixteenth1916 & \Who{C2}{Brice Dudout} & 1916 & Chinese & CT & complete & 168 \\
thirteenth1916 & \Who{C2}{Brice Dudout} & 1916 & Chinese & CT & complete & 100 \\
thirtyeighth1916 & \Who{C2}{Brice Dudout} & 1916 & Chinese & CT & complete & 40 \\
thirtyfifth1916 & \Who{C2}{Brice Dudout} & 1916 & Chinese & CT & complete & 130 \\
wenmi1916g0 & \Who{C2}{Brice Dudout} & 1916 & Chinese & CT & complete & 351 \\
wenmi1916y0 & \Who{C2}{Brice Dudout} & 1916 & Chinese & CT & complete & 456 \\
ylchang1916 & \Who{C2}{Brice Dudout} & 1916 & Chinese & CT & complete & 266 \\
bernstorff1917 & \Who{C3}{WAINUT} & 1917 & German & EPT & partial & 577 \\
visser1917 & \Who{C3}{WAINUT} & 1917 & Spanish & NS & none & -- \\
auckland1918 & \Who{C2}{Brice Dudout} & 1918 & English & CT & complete & 399 \\
cpa1920 & \Who{C3}{WAINUT} & 1920 & English & APT & partial & 439 \\
tokyo1932 & \Who{C2}{Brice Dudout} & 1932 & English & DEC & partial & -- \\
moscowmeteo1940 & \Who{C2}{Brice Dudout} & 1940 & International & DEC & complete & 17 \\
ts1941pz6 & \Who{C3}{WAINUT} & 1941 & German & NS & none & 683 \\
bertelot & \Who{C1}{George Lasry} & -- & French & APT & partial & 1646 \\
clair1161 & \Who{C1}{George Lasry} & -- & French & APT & partial & 3149 \\
clair460bar & \Who{C1}{George Lasry} & -- & French & APT & partial & 1218 \\
clair460num & \Who{C1}{George Lasry} & -- & French & APT & partial & 2080 \\
decode4282 & \Who{C1}{George Lasry} & -- & Latin & CT & complete & 1084 \\
dubellay65 & \Who{C1}{George Lasry} & -- & French & CT & partial & 851 \\
fr14765p351 & \Who{C1}{George Lasry} & -- & French & CT & partial & 936 \\
fr3321dc & \Who{C1}{George Lasry} & -- & French & APT & partial & 49 \\
fr3575f56 & \Who{C1}{George Lasry} & -- & French & CT & complete & 782 \\
lerouge & \Who{C1}{George Lasry} & -- & French & APT & partial & 623 \\
siena13 & \Who{C1}{George Lasry} & -- & Italian & APT & complete & 553 \\
siena5 & \Who{C1}{George Lasry} & -- & Italian & CT & complete & 165 \\
sweden4306 & \Who{C1}{George Lasry} & -- & Dutch & CT & partial & 461 \\
villeroy258 & \Who{C1}{George Lasry} & -- & French & APT & complete & 1164 \\
\end{longtable}}


\end{document}
